\documentclass[11pt]{article}
\usepackage[margin=27mm]{geometry}
\usepackage[T1]{fontenc}
\usepackage{lmodern}
\usepackage{amsmath,amssymb,amsthm,mathtools,mathrsfs}
\usepackage{booktabs,array,graphicx,microtype}
\usepackage[numbers,sort&compress]{natbib}
\usepackage[colorlinks=true,allcolors=blue]{hyperref}
\usepackage{xurl}
\usepackage{doi}
\newtheorem{theorem}{Theorem}[section]
\newtheorem{proposition}[theorem]{Proposition}
\newtheorem{lemma}[theorem]{Lemma}
\newtheorem{corollary}[theorem]{Corollary}
\theoremstyle{definition}
\newtheorem{conjecture}[theorem]{Conjecture}
\newtheorem{definition}[theorem]{Definition}
\theoremstyle{remark}
\newtheorem{remark}[theorem]{Remark}
\newcommand{\CC}{\mathbb C}
\newcommand{\QQ}{\mathbb Q}
\newcommand{\ZZ}{\mathbb Z}
\newcommand{\FF}{\mathbb F}
\newcommand{\remA}{\operatorname{rem}_{A}}
\newcommand{\len}{\operatorname{length}}
\newcommand{\ct}{\operatorname{ct}}
\newcommand{\Xp}{\mathcal X_p}
\newcommand{\Wclass}{\mathcal W}
\newcommand{\dd}{\,\mathrm d}
\newcommand{\fall}[2]{#1^{\underline{#2}}}
\DeclareMathOperator{\Res}{Res}
\numberwithin{equation}{section}
\setlength{\emergencystretch}{2em}
\renewcommand{\topfraction}{0.9}
\renewcommand{\bottomfraction}{0.8}
\renewcommand{\textfraction}{0.1}
\renewcommand{\floatpagefraction}{0.8}
\title{A binomial count of meromorphic travelling waves}
% Author metadata.
\input{authors}
\date{}
\begin{document}
\maketitle

\begin{abstract}
We study meromorphic travelling waves of
\(u_t+uu_x+\sum_{j=1}^{p}b_j\partial_x^{j+1}u=0\), a
constant-coefficient family containing KdV--Burgers,
Kuramoto--Sivashinsky and Kawahara. After integration and
normalization, the profiles satisfy \(P(D)v+v^2/2=0\), where
\(D=\mathrm d/\mathrm dz\) and \(P\) is monic of degree \(p\ge2\) with
complex coefficients.

For evolution equations of even order (\(p\) odd), and for purely
dispersive equations (only odd-order derivatives), we give a complete
algebraic description of all nonconstant globally meromorphic
travelling waves: they are rational, rational in an exponential, or
elliptic, with one pole per period.

One-pole waves of these three kinds are exceptional: for \(p\ge3\) a
generic equation of the family has none. We count the exceptions as
pairs \((P,v)\), over the whole family: there are
\(\binom{2p-1}{p-1}\) with \(v\) rational in an exponential of given
rate, and \(2\binom{2p-2}{p-2}\) with \(v\) elliptic on a given generic
lattice. We also count the solitary pulses of the purely dispersive
subfamily and the fronts of the purely dissipative one (only even-order
derivatives in the linear part).

When \(2p-1\) is prime, these statements are proved and all pairs are
distinct. We conjecture them in every order.
\end{abstract}

\noindent\textbf{Keywords.}
Meromorphic travelling waves; dispersive and dissipative equations;
elliptic functions; polynomial convolution; algebraic multiplicity.

\section{Introduction}\label{sec:intro}

Exact travelling waves of nonlinear evolution equations are usually
found one equation at a time, and usually only for special values of
the coefficients. The generalized Kuramoto--Sivashinsky equation
\(u_t+uu_x+u_{xx}+\sigma u_{xxx}+u_{xxxx}=0\) is a classical example.
It has nonconstant travelling waves that are polynomials in a
hyperbolic tangent for exactly four values of the dispersion
coefficient, up to its sign: \(\sigma=0\), \(4\), \(12/\sqrt{47}\) and
\(16/\sqrt{73}\). Kuramoto and Tsuzuki found the wave of the
dispersionless equation \citep{KuramotoTsuzuki1976}, Kudryashov the
others \citep{Kudryashov1988,Kudryashov1990}, and Eremenko proved
that, together with Kudryashov's elliptic solution and a single rational
one, they exhaust the nonconstant meromorphic travelling waves
\citep{Eremenko2006}.
The classical tables list six cases at these four values
\citep{KudryashovZargaryan1996,ConteMusette2009}. They do not
distinguish a wave from its mirror image, which solves the equation
with the opposite sign of \(\sigma\). If the two are counted
separately, the two cases with \(\sigma=0\) are their own mirror images
and the other four are not, so the list has \(2+2\cdot4=10\) members.
This paper shows that ten is \(\binom52\): the classical list is the
case \(p=3\) of a binomial enumeration, proved for infinitely many
\(p\) and conjectured for all.

We consider the evolution equations
\[
 u_t+uu_x+\sum_{j=1}^{p}b_j\partial_x^{j+1}u=0,\qquad b_p\ne0,
\]
which contain the KdV--Burgers (\(p=2\)), Kuramoto--Sivashinsky
(\(p=3\)), Kawahara and Benney--Lin (\(p=4\)) and Nikolaevskiy
(\(p=5\)) equations; Table~\ref{tab:named} lists the named members with
\(p\le6\).
A travelling wave \(u=w(x-ct)\) satisfies, after one integration,
\(\sum_jb_jw^{(j)}-cw+w^2/2+K=0\). Shifting \(w\) by a constant
solution and normalizing the leading coefficient gives
\begin{equation}\label{eq:main}
 E_P(v):=P(D)v+\frac12v^2=0,\qquad
 P(\rho)=\rho^p+\sum_{j=0}^{p-1}a_j\rho^j,\qquad p\ge2,
\end{equation}
with \(D=\mathrm d/\mathrm dz\) and complex coefficients. The
coefficients \(a_1,\ldots,a_{p-1}\) are those of the evolution
equation, up to a rescaling of \(x\), and the speed, the integration
constant and the background are absorbed into \(a_0\)
(Section~\ref{sec:setting}).
Throughout, \(p\) is the order of the profile equation \eqref{eq:main},
which is one less than the order of the evolution equation:
Kuramoto--Sivashinsky is a fourth-order equation and has \(p=3\),
Kawahara has \(p=4\), and the seventh-order KdV equation has \(p=6\).
Phrases such as ``order three'' or ``in every order'' always refer to
\(p\).

\paragraph{The waves.}
A travelling wave is \emph{meromorphic} if its profile is a
single-valued meromorphic function of the complex variable \(z\). The
familiar exact solutions are of this kind: a polynomial in
\(\tanh\) or \(\operatorname{sech}\) is a rational function of an
exponential \(e^{\kappa z}\), a cnoidal wave is an elliptic function,
and rational functions of \(z\) also occur. These are the three kinds
of wave in this paper. We consider those with \emph{one pole per
period}: all the poles of the profile are translates of a single pole
by its periods. For a wave of \(\tanh\) type the period in question is
the imaginary period \(2\pi i/\kappa\) of \(e^{\kappa z}\), and the
rate \(\kappa\) measures the width of the pulse or front. For many
equations of the family, described below, these are all the meromorphic
travelling waves that exist.

\paragraph{The question.}
For \(p\ge3\) such waves are exceptional: an equation of the family has
one only if its coefficients satisfy algebraic relations, as the four
values of \(\sigma\) illustrate. This suggests the right question. We
do not ask how many waves a given equation has, but how many pairs of
an equation and a wave there are in the whole family.

To make this question finite, two continuous freedoms must be removed.
A translate of a wave is again a wave, so we place the pole at a fixed
point. Rescaling \(x\) changes the width of a wave and the coefficients
of its equation together, so we fix the width: for a wave that is
rational in an exponential we prescribe the rate \(\kappa=1\), and for
an elliptic wave we prescribe both of its periods, that is, its period
lattice. The coefficients of the equation remain free. We therefore
count pairs \((P,v)\), where \(P\) ranges over all monic polynomials of
degree \(p\) and \(v\) solves \eqref{eq:main}:
\begin{itemize}
\item a \emph{rational-exponential pair} has \(v\) rational in
\(e^z\), with one pole modulo \(2\pi i\ZZ\), placed at \(i\pi\), and
\(v\to0\) as \(e^z\to0\);
\item an \emph{elliptic pair} has \(v\) elliptic on a given lattice,
with its only pole at zero.
\end{itemize}
One more convention explains a factor two below. The profile equation
has two constant solutions, and a wave can be measured from either of
them; this changes \(P\) but not the wave (Section~\ref{sec:setting}).
A rational-exponential wave is measured from its limit as \(e^z\to0\).
An elliptic wave has no such limit, and both choices are counted, so
every elliptic wave appears in two pairs.

\paragraph{The answer.}
The numbers of pairs are
\[
 \begin{aligned}
 N_p&=\binom{2p-1}{p-1}
     &&\text{rational-exponential pairs},\\
 M_p&=2\binom{2p-2}{p-2}
     &&\text{elliptic pairs on a generic lattice}.
 \end{aligned}
\]
We call \(2p-1\) the \emph{index} of the order \(p\). Our main result
establishes these numbers whenever the index is a prime number:
\begin{itemize}
\item there are exactly \(N_p\) rational-exponential pairs, and they
are all distinct (Theorem~\ref{thm:prime});
\item exactly \(S_p=\binom{2p-2}{p-2}\) of them are \emph{pulses},
which return to the same constant at both ends, and the others are
\emph{fronts}, which connect two different constants
(Proposition~\ref{cor:endpoints});
\item a generic lattice carries exactly \(M_p=2S_p\) elliptic pairs,
and they are all distinct (Theorem~\ref{thm:elliptic-count}).
\end{itemize}
The index is prime for \(p=2,3,4,6,7,9,10,12,\ldots\). At the three
remaining orders up to twelve, \(p=5\), \(8\) and \(11\), the same
statements are established by exact computer calculations
(Corollary~\ref{cor:twelve}). We call the orders with a prime index,
together with all \(p\le12\), the \emph{covered} orders; they include
every named equation of Table~\ref{tab:named}. We conjecture the same
statements in every order (Conjecture~\ref{conj:count}).

\begin{table}[htbp]
\centering
\caption{The named equations of the family with \(p\le6\); here
\(u_{kx}=\partial_x^ku\). Each name stands for a family with free
coefficients, and the numbers count distinct pairs across that family:
rational-exponential pairs at unit rate, and elliptic pairs on a
generic given lattice, where every elliptic wave appears twice. An
equation with only odd-order, or only even-order, linear terms belongs
to a pure subfamily, counted by \(F_p\) and \(2F_p\) (see below); the
full families are counted by \(N_p\) and \(M_p\) and contain the pure
subfamily of the same \(p\), so the three pairs of KdV--Burgers include
the KdV soliton. Except in the two rows marked \(\dagger\), these pairs
are all the nonconstant meromorphic travelling waves of the equations,
apart from a rational wave when all lower coefficients vanish.}
\label{tab:named}
\vspace{4pt}
\begin{tabular}{@{}llcrr@{}}
\toprule
Equation&Linear terms&\(p\)&Rational-exponential pairs&Elliptic pairs\\
\midrule
KdV&\(u_{3x}\)&2&\(F_2=1\)&\(2F_2=2\)\\
KdV--Burgers\(^\dagger\)&\(u_{2x},u_{3x}\)&2&\(N_2=3\)&\(M_2=2\)\\
Kuramoto--Sivashinsky (KS)&\(u_{2x},u_{4x}\)&3&\(F_3=2\)&\(0\)\\
KS with dispersion&\(u_{2x},u_{3x},u_{4x}\)&3&\(N_3=10\)&\(M_3=8\)\\
Kawahara&\(u_{3x},u_{5x}\)&4&\(F_4=3\)&\(2F_4=6\)\\
Benney--Lin\(^\dagger\)&\(u_{2x},\ldots,u_{5x}\)&4&\(N_4=35\)&\(M_4=30\)\\
Nikolaevskiy&\(u_{2x},u_{4x},u_{6x}\)&5&\(F_5=6\)&\(0\)\\
Nikolaevskiy with dispersion&\(u_{2x},\ldots,u_{6x}\)&5&\(N_5=126\)&\(M_5=112\)\\
seventh-order KdV&\(u_{3x},u_{5x},u_{7x}\)&6&\(F_6=10\)&\(2F_6=20\)\\
\bottomrule
\end{tabular}
\end{table}

Table~\ref{tab:named} lists the counts for the named equations of the
family with \(p\le6\): the Benney--Lin equation
\citep{Benney1966,Lin1974}, the Kuramoto--Sivashinsky equation with
dispersion, also called the Benney or KdV--KS equation
\citep{KawaharaToh1988}, the Nikolaevskiy equation
\citep{Nikolaevskii1989} and its dispersive version
\citep{KudryashovMigita2007,SimbawaMatthewsCox2010}, the long-wave
reductions of \citet{KudryashovSinelshchikovChernyavsky2008}, and the
equations already mentioned. The small entries are familiar. The three
pairs of KdV--Burgers are the KdV soliton and the KdV--Burgers front,
read in its two directions. The two elliptic pairs of KdV are the
cnoidal wave on its two backgrounds. The eight elliptic pairs at
\(p=3\) are Kudryashov's elliptic solution: on a given lattice it
occurs for four values of the dispersion coefficient, each time on two
backgrounds (Section~\ref{sec:ks}).

Four remarks make precise what the numbers mean.
The counts are over \(\CC\) and range over all monic \(P\). Many of the
pairs have real coefficients and give smooth bounded waves of real
equations: eight of the ten at \(p=3\), and \(21\) of the \(35\) at
\(p=4\); Section~\ref{sec:examples} identifies them.
The word \emph{distinct} carries content: the pairs are the solutions
of a polynomial system with as many equations as unknowns, and a
solution of such a system can be multiple, like a multiple root of a
polynomial. On special lattices elliptic pairs do merge in this way,
and Section~\ref{sec:ks} determines where this happens at order three.
Rational waves are not counted: at the covered orders the only one that
vanishes at infinity is \(v=c_\ast z^{-p}\), for \(P=D^p\).
Finally, the waves of low order are classical, and
Section~\ref{sec:examples} gives the references; what is new is that
the lists are complete, and how long they are. We do not discuss
stability, and we treat only the quadratic nonlinearity \(uu_x\).

\paragraph{Every wave has a label.}
Here is where the binomial coefficient \(N_p\) comes from. Put
\(Q=1/(1+e^{-z})=(1+\tanh(z/2))/2\). Every rational-exponential pair
has a profile of the form \(v_A=s_pA(D)Q\), where \(s_p\) is a fixed
scale and \(A\) is a monic polynomial of degree \(p-1\); this is a
polynomial in \(\tanh(z/2)\) written in a convenient basis. For
\(\operatorname{Re}z<0\) we have \(Q=e^z-e^{2z}+e^{3z}-\cdots\), and
\(D\) multiplies \(e^{nz}\) by \(n\). Hence the coefficient of
\(e^{nz}\) is \(\pm s_pA(n)\) in \(v_A\) and \(\pm s_pP(n)A(n)\) in
\(P(D)v_A\), while in \(v_A^2\) it is \(\pm s_p^2S_A(n)\) with the
convolution sum
\[
 S_A(n)=\sum_{j=1}^{n-1}A(j)A(n-j).
\]
These sums are the values of a polynomial \(S_A\) of degree \(2p-1\),
and \eqref{eq:main} holds exactly when \(2P(n)A(n)=s_pS_A(n)\) for all
\(n\). So the pair exists if and only if the polynomial \(A\) divides
the polynomial \(S_A\), and then \(P\) is the quotient, up to the
constant factor (Proposition~\ref{prop:divisibility}).

The smallest case shows what happens. For \(p=2\) and \(A=\rho+a\),
\[
 S_A(\rho)=(\rho-1)\Bigl(\frac{\rho(\rho+1)}6+a\rho+a^2\Bigr),
\]
which is divisible by \(\rho+a\) exactly when \(a(a-1)(a+1)=0\). The
solution \(a=0\) is the KdV soliton, and \(a=\pm1\) are the
KdV--Burgers front read in its two directions. Now reduce modulo
\(3=2p-1\). Whatever the value of \(a\),
\[
 3S_A(\rho)\equiv-(\rho^3-\rho)=-\rho(\rho-1)(\rho+1)\pmod3,
\]
so modulo \(3\) the divisor \(A\) must be one of the three factors
\(\rho\), \(\rho-1\), \(\rho+1\). Each of them comes from exactly one
solution, and the root of \(A\) modulo \(3\) serves as a label: \(\{0\}\)
for the soliton, \(\{2\}\) and \(\{1\}\) for the two fronts.

The same happens whenever \(\ell=2p-1\) is prime. The polynomial
\(S_A\) is a combination of the binomial coefficients
\(\binom{\rho-1}{k}\), and only the last of them, with \(k=\ell\), has
\(\ell\) in its denominator. For this reason \(\ell S_A\) reduces
modulo \(\ell\) to \(\pm(\rho^\ell-\rho)\), whatever \(A\) is, and the
condition becomes that \(A\) divides \(\rho^\ell-\rho\) modulo \(\ell\).
Now \(\rho^\ell-\rho\) is the product of \(\rho-s\) over all residues
\(s\) modulo \(\ell\). Its monic divisors of degree \(p-1\) are
therefore given by the \((p-1)\)-element subsets \(S\) of the integers
modulo \(\ell\), and there are \(\binom{\ell}{p-1}=N_p\) of them. We
prove that every pair can be reduced modulo \(\ell\), although its
coefficients are in general algebraic numbers and not integers, and
that every divisor comes from exactly one pair. For instance
\(A=\rho^2+1/11\), the Kuramoto--Tsuzuki wave, reduces to
\(\rho^2+1=(\rho-2)(\rho-3)\) modulo \(5\), and has the label
\(\{2,3\}\).

So each pair has a label \(S\), and the labels carry information.
The pair is a pulse exactly when \(S\) contains \(0\), and the mirror
image of the pair with label \(S\) has label \(-S\). At \(p=3\) the ten
classical waves correspond to the ten two-element subsets of
\(\{0,1,2,3,4\}\) (Table~\ref{tab:ks}), and the four pulses are the
subsets that contain \(0\).

The elliptic count comes from the pulses. When one period of the
lattice tends to infinity, an elliptic pair becomes a pulse, measured
from one of its two constant backgrounds. We show that every pulse, on
either background, arises in this way exactly once, and this gives
\(M_p=2S_p\).

\paragraph{Purely dispersive and purely dissipative equations.}
In the evolution equation, derivatives of odd order are dispersive
terms and derivatives of even order are dissipative terms. We call an
equation \emph{purely dispersive} if it contains only odd-order
derivatives, as the KdV and Kawahara equations do; then \(p\) is even
and \(P\) is an even polynomial. We call it \emph{purely dissipative}
if its linear part contains only even-order derivatives, as for the
Kuramoto--Sivashinsky equation without dispersion; then \(p\) is odd
and \(P-P(0)\) is an odd polynomial. All other equations, such as
KdV--Burgers, are \emph{mixed}.

The two pure subfamilies are singled out by reflection. Reading a wave
in the opposite direction, \(z\mapsto-z\), gives a wave of the
mirror-image equation, in which the terms of one parity have changed
sign. For a front the zero limit is then at the wrong end, so one also
changes the background; the result is again a pair, the \emph{mirror
image} of the first (Section~\ref{sec:reflection}). The pairs that
coincide with their mirror image are exactly those whose equation is
purely dispersive, for even \(p\), or purely dissipative, for odd \(p\)
(Theorem~\ref{thm:symmetric}). In terms of labels they are the pairs
with \(S=-S\), and the number of such subsets is
\[
 F_p=\binom{p-1}{\lfloor(p-1)/2\rfloor}.
\]
At the covered orders this is the number of rational-exponential pairs
in the pure subfamily of order \(p\), and the purely dispersive
equations also account for \(2F_p\) of the elliptic pairs on a generic
lattice (Corollary~\ref{cor:elliptic-symmetric}).

The two subfamilies carry different kinds of rational-exponential
waves. Those of a purely dispersive equation are all pulses, at every
even \(p\). Those of a purely dissipative equation are all fronts, at
the covered orders; for waves with real coefficients this holds at
every odd \(p\). For example, the Kawahara family has three
rational-exponential pairs and six elliptic pairs on a generic lattice,
which are \(F_4\) and \(2F_4\); the seventh-order KdV family has ten
pairs, two of them real; and the six pairs of the purely dissipative
family with \(p=5\) contain the four kink solutions of
\citet{KudryashovMigita2007} (Section~\ref{sec:examples}).

\paragraph{When these are all the meromorphic waves.}
In two cases the one-pole waves of the three kinds are all the
nonconstant meromorphic travelling waves there are
(Theorem~\ref{thm:complete}). The first case is \(p\) odd, that is,
every evolution equation of even order. The second is the purely
dispersive case. The proof rests on Eremenko's theorem
\citep{Eremenko2006}: an autonomous algebraic differential equation
has only these three kinds of meromorphic solutions if it has a single
monomial of top degree, here \(v^2/2\), and finitely many Laurent
expansions at a pole. For odd \(p\) the Laurent expansion at a pole is
unique. For a purely dispersive equation one Laurent coefficient is
free, a resonance; but the equation is invariant under \(z\mapsto-z\)
and conserves an energy, and fixing the energy fixes that coefficient.
In both cases every nonconstant meromorphic travelling wave has one
pole per period, and for a given operator \(P\), which records the
equation and the speed, there is at most one such profile up to
translation when \(p\) is odd, and at most one on each energy level
when the equation is purely dispersive.

Mixed operators with \(p\) even, such as that of KdV--Burgers, are not
covered by this theorem, and cannot be: Section~\ref{sec:open} recalls
an example with \(p=2\) that has other meromorphic waves. For these
operators we still count the one-pole waves of the three kinds, without
claiming that there are no others; they are the rows marked
\(\dagger\) in Table~\ref{tab:named}. Section~\ref{sec:systems} writes
each of the three kinds as a finite polynomial system, following the
principal-part construction of
\citet{DeminaKudryashov2011,DeminaKudryashovElliptic2011}.

\paragraph{A single equation.}
A count across the family is not a count for each individual equation.
To recover the waves of a given equation, we restore the length scale
and retain only the pairs whose coefficients match it. Changing the
scale traces a curve in the space of coefficients, and at the covered
orders the equations with a one-pole rational-exponential wave are
those on finitely many such curves, at most one for each pair and its
mirror image. This is
the precise sense in which the waves are exceptional, and it explains
the four values of \(\sigma\). It refers to the full family: a pure
subfamily with few coefficients, such as the KdV or Kawahara equations,
can lie entirely on these curves, and at \(p=2\) every equation does,
since every KdV--Burgers equation has its front. If at least one
lower-derivative coefficient is nonzero, an equation has only finitely
many wave shapes after translations and Galilean shifts are identified;
the scale-invariant case is different, and the KdV equation, for
example, retains its family of solitons of all widths.
Section~\ref{sec:fixed} makes this precise and recovers the four
Kuramoto--Sivashinsky curves.

For equations of even order the classification gives more. A generic
such equation has no nonconstant meromorphic travelling wave at all;
this is proved for \(3\le p\le573\) (Corollary~\ref{cor:generic}). It may well have travelling
waves; but none of them is meromorphic, and in particular none is given
by a formula of the kinds above. Here generic means outside a proper
algebraic subset of the coefficients.

Section~\ref{sec:fixed} also records two consequences for the way exact
travelling waves are found in practice. Many other equations, with
mixed derivatives or in several space variables, have the same profile
equation, so that the counts apply to them. And the expansion methods
by which such waves are usually sought, in powers of \(\tanh\) or of a
Jacobi function, can only produce waves that are counted here; for a
named equation the counts are a checklist for their output.

\paragraph{Other orders.}
The proofs use the prime \(\ell\) in two places, and these are the two
facts that are not yet proved in every order. The first is
\emph{rigidity}: the only rational solution of \eqref{eq:main} with
one pole that vanishes at infinity is \(c_\ast z^{-p}\), with
\(P=D^p\). Equivalently, the only monic polynomial \(A\) of degree
\(p-1\) that divides its continuous self-convolution
\(\int_0^\rho A(t)A(\rho-t)\dd t\) is \(\rho^{p-1}\). As the
coefficients of the equation vary, a pair could disappear only by
moving off to infinity, and rigidity is the statement that this does
not happen; it gives all the counts when pairs are counted with
multiplicity. It is proved when \(2p-1\) is a prime or a prime power,
and established by exact computations for \(p\le573\). The second fact
is that every rational-exponential pair is simple, which makes the
pairs distinct. It is proved when \(2p-1\) is prime and, by exact
computations, at \(p=5\), \(8\) and \(11\). Both are conjectured in
general. Section~\ref{sec:further} states what follows from each of
them, and a companion note in the source package contains the proofs of
those statements.

\paragraph{Relation to earlier work.}
The classification theorem is a uniform version of results that exist
in pieces. The case of odd \(p\) is contained in
\citet[Theorem~1]{YuanLiQi2014}, and the energy device for even \(P\) is
that of \citet[Lemmas~3--4]{EremenkoLiaoNg2009}, which has been applied
to the Kawahara, generalized Bretherton and seventh-order KdV
equations; Section~\ref{sec:completeness} gives the detailed
attributions. The polynomial systems are instances of the
principal-part construction of
\citet{DeminaKudryashov2011,DeminaKudryashovElliptic2011}.
The general enumeration has no precedent that we know of; its special
cases at \(p=3\) and in the Kawahara family are classical, as
the examples make explicit.

\paragraph{Organization.}
Section~\ref{sec:setting} contains the normalizations, and
Sections~\ref{sec:completeness} and~\ref{sec:systems} the
classification and the three polynomial systems.
Section~\ref{sec:counting} proves the count of rational-exponential
pairs at a prime index and reads the labels, and
Section~\ref{sec:elliptic-counting} proves the elliptic count.
Section~\ref{sec:further} records what is known at the other orders,
and Section~\ref{sec:fixed} reads the count for a single
equation. Section~\ref{sec:examples} works out the lowest orders and
the two pure subfamilies, including the lattices on which elliptic
pairs merge. Section~\ref{sec:open} collects the open questions, and
Appendix~\ref{app:verification} describes the exact verification and
the ancillary material.

\section{Equations, waves and what is counted}\label{sec:setting}

This section fixes the normalizations and records the form of the waves
that are counted.

\subsection{From the evolution equation to the profile equation}

For a travelling wave \(u=w(\xi)\), \(\xi=x-ct\), one integration turns
\eqref{eq:pde} into \eqref{eq:integrated}:
\begin{align}
 &u_t+uu_x+\sum_{j=1}^{p}b_j\partial_x^{j+1}u=0,\qquad b_p\ne0,
 \label{eq:pde}\\
 &\sum_{j=1}^{p}b_jw^{(j)}-cw+\frac12w^2+K=0,\label{eq:integrated}
\end{align}
with a constant \(K\). Let \(\kappa\ne0\), and let \(w_\ast\) be a
constant solution of \eqref{eq:integrated}, a root of \(w^2/2-cw+K\).
The substitution \eqref{eq:scaling} turns \eqref{eq:integrated} into the
profile equation \eqref{eq:main} with coefficients
\eqref{eq:coefficients}:
\begin{align}
 w(\xi)&=w_\ast+b_p\kappa^p\,v(\kappa\xi),\label{eq:scaling}\\
 a_j&=\frac{b_j}{b_p}\kappa^{j-p}\quad(1\le j\le p-1),\qquad
 a_0=\frac{w_\ast-c}{b_p\kappa^p}.\label{eq:coefficients}
\end{align}
Thus \(a_1,\ldots,a_{p-1}\) record the equation, up to the scale
\(\kappa\), and \(a_0\) records the speed of the wave relative to the
background \(w_\ast\). The scale is free at this stage. The period of
the wave will fix it: if \(v\) is rational in \(e^{z}\), then \(w\) is
rational in \(e^{\kappa\xi}\), so that \(\kappa\) is the rate of the wave.

The constant solutions of \eqref{eq:main} are \(v=0\) and \(v=-2a_0\).
They correspond to the two possible choices of the background
\(w_\ast\). Changing that choice replaces the pair \((P,v)\) by
\begin{equation}\label{eq:shift}
 \bigl(P-2P(0),\ v+2P(0)\bigr),
\end{equation}
which again satisfies \eqref{eq:main}, as a direct substitution shows.
In the evolution equation the change of background is a Galilean shift
\(u\mapsto u+k\), \(x\mapsto x-kt\); it changes the speed relative to
the background, hence the sign of \(a_0\). The two pairs describe the
same solution \(u(x,t)\) of the same evolution equation. They coincide
exactly when \(a_0=0\). For example, the KdV soliton gives the pair
\(P=\rho^2-1\), \(v=3\operatorname{sech}^2(z/2)\). Measured from the
other background it is \(P=\rho^2+1\) and \(v-2\), which does not tend
to zero and is therefore not a rational-exponential pair.

\subsection{Dispersive, dissipative and mixed equations}

In \eqref{eq:pde} the term \(b_j\partial_x^{j+1}u\) is dispersive for
even \(j\) and dissipative for odd \(j\): for real coefficients it
changes the phase, respectively the amplitude, of a Fourier mode (with
the signs of \eqref{eq:pde}, a positive coefficient of \(u_{xx}\) is
gain). The equation is \emph{purely dispersive} if \(b_j=0\) for odd
\(j\). Then \(p\) is even, \(P\) is an even polynomial, and the equation
is invariant under \((x,t)\mapsto(-x,-t)\). The equation is \emph{purely
dissipative} if \(b_j=0\) for even \(j\). Then \(p\) is odd, \(P-P(0)\)
is an odd polynomial, and the equation is invariant under
\((x,u)\mapsto(-x,-u)\); the Nikolaevskiy equation (\(p=5\)) is an
example, see \citet{SimbawaMatthewsCox2010}. All other equations are
\emph{mixed}.

The two pure subfamilies are described by one condition. We call a monic
polynomial \(P\) of degree \(p\) \emph{reflection-symmetric} if
\(P-P(0)\) has the parity of \(\rho^p\):
\begin{equation}\label{eq:symmetric-class}
 P(-\rho)=P(\rho)\quad(p\text{ even}),\qquad
 P(\rho)+P(-\rho)=2P(0)\quad(p\text{ odd}).
\end{equation}
In the travelling-wave variable both symmetries are the reflection
\(z\mapsto-z\). If \(P\) is even and \(v(z)\) solves \eqref{eq:main}, so
does \(v(-z)\). If \(p\) is odd and \(P-P(0)\) is odd, so does
\(-v(-z)-2P(0)\); this is the reflection combined with the change of
background \eqref{eq:shift}.

\subsection{Poles, and the form of one-pole waves}\label{sec:systems}

To count waves we need to know what they look like. Three facts suffice:
a pole has order \(p\) and a fixed leading coefficient; with one pole
per period, a wave rational in \(e^z\) is a polynomial of degree \(p\)
in the logistic function \(Q\), and an elliptic wave is a combination of
\(\wp\) and its derivatives.

Throughout, a meromorphic function is single-valued and meromorphic on
the whole complex plane. The order of a pole of a solution of
\eqref{eq:main} is forced by dominant balance. A pole of order \(r\)
makes \(P(D)v\) of order \(r+p\) and \(v^2\) of order \(2r\); the two
can cancel only if \(r=p\). Comparing the leading coefficients gives, at
a pole \(z_0\),
\begin{equation}\label{eq:leadingpole}
 v(z)=\frac{c_\ast}{(z-z_0)^{p}}+O\bigl((z-z_0)^{1-p}\bigr),\qquad
 c_\ast=(-1)^{p-1}(p-1)!\,s_p,\qquad s_p=\frac{2(2p-1)!}{((p-1)!)^2}.
\end{equation}
This holds for every \(P\). The scale \(s_p\) of the normal forms of
Sections~\ref{sec:counting} and~\ref{sec:elliptic-counting} has the
first values \(s_2=12\), \(s_3=60\), \(s_4=280\), \(s_6=5544\), so that
\(c_\ast=-12,\,120,\,-1680\) for \(p=2,3,4\).

The meromorphic travelling waves in this paper are of three kinds:
rational in \(z\), rational in an exponential \(e^{\kappa z}\), and
elliptic; Theorem~\ref{thm:complete} says when there are no others. The
period group of a wave is its \emph{full} group of periods: trivial in
the first case, generated by one period in the second, a lattice in the
third. In the second case we choose \(\kappa\) so that
\(2\pi i/\kappa\) generates this group. The wave has \emph{one pole per
period} if all its poles are translates of a single pole by the periods.
A rational wave then has a single pole, and an elliptic wave one in each
period parallelogram.

Let \(v=R(t)\), \(t=e^{\kappa z}\), be a wave rational in an
exponential, with one pole per period. The rational function \(R\) has
finite limits at both \(t=0\) and \(t=\infty\). Indeed,
\(P(D)=P(\kappa t\partial_t)\) multiplies each monomial \(t^m\) by a
constant, whereas squaring doubles an extreme exponent: if \(R\) had a
pole at either end, \(R^2\) would contain an extreme Laurent monomial
that \(P(\kappa t\partial_t)R\) cannot cancel. Both limits are constant
solutions of \eqref{eq:main}. After a translation, the only pole of
\(R\) in \(0<|t|<\infty\) is at \(t=-1\), of order \(p\). For
\(\kappa=1\) such a function is a polynomial of degree \(p\) in the
logistic function \(Q=1/(1+e^{-z})=(1+\tanh(z/2))/2\), with
\(DQ=Q(1-Q)\), which has a simple pole with residue one at \(z=i\pi\)
and the limits \(0\) and \(1\). By \eqref{eq:leadingpole} the
coefficient of \(Q^p\) is \(c_\ast\); the KdV soliton above is
\(12Q(1-Q)\). Polynomial logistic representations and their use in
constructing travelling waves are developed by
\citet{Kudryashov2012,Kudryashov2013,Kudryashov2015}, and applied to
equations of the family up to \(p=6\) by
\citet{RyabovSinelshchikovKochanov2011}.

Such a wave is a \emph{pulse} if its two limits coincide and a
\emph{front} otherwise. For a real rate and a real profile these are the
solitary waves and the kinks of the evolution equation, and they are
smooth and bounded on the real axis when the pole lies off it. In
\eqref{eq:scaling} we take \(w_\ast\) to be the limit as
\(e^{\kappa z}\to0\). Then \(v\to0\) at that end, and the other limit
is \(0\) for a pulse and \(-2a_0\) for a front.

An elliptic wave whose only pole modulo its periods is at zero is a
constant plus a linear combination of \(\wp,\wp',\ldots,\wp^{(p-2)}\),
in which \(\wp^{(p-2)}\) has the coefficient \(-s_p\)
(Section~\ref{sec:elliptic-counting}). Rational waves are not counted,
but they govern what happens at infinity in both counts. A rational
solution has no nonconstant polynomial part at infinity: such a part
would have its degree doubled by \(v^2\), while the linear differential
expression cannot reach that degree. A rational wave whose only pole is
at zero is therefore
\begin{equation}\label{eq:rational-system}
 v(z)=c_\ast z^{-p}+\bigl(\text{a polynomial of degree at most }p-1
 \text{ in }z^{-1}\bigr).
\end{equation}
The simplest one is \(v=c_\ast z^{-p}\), for \(P=D^p\); it is given for
every order by \citet{Kudryashov2004}.

Conversely, a function of one of these forms that satisfies
\eqref{eq:main} is a nonconstant wave with one pole per period, because
\(c_\ast\ne0\). Substitution turns \eqref{eq:main} into finitely many
polynomial equations in the coefficients of the form and of \(P\),
without any genericity assumption. They are instances of the
principal-part construction of
\citet{DeminaKudryashov2011,DeminaKudryashovElliptic2011}; for a fixed
operator it reduces to the coefficient comparison of
\citet{DeminaKudryashov2014}, which we do not need.

\subsection{What is counted}

We count, over \(\CC\), pairs \((P,v)\) in which \(P\) ranges over
\emph{all} monic polynomials of degree \(p\) and \(v\) solves
\eqref{eq:main}. In a \emph{rational-exponential pair}, \(v\) is
rational in \(e^z\), its poles are the points of \(i\pi+2\pi i\ZZ\), and
\(v\to0\) as \(e^z\to0\). In an \emph{elliptic pair} on a lattice
\(\Lambda\), \(v\) is elliptic with period lattice \(\Lambda\), including
its scale, and its poles are the points of \(\Lambda\). Fixing the
periods removes the scale \(\kappa\) in \eqref{eq:scaling}, and placing
the pole removes translations. A wave and its mirror image are counted
as two pairs when they are distinct (Section~\ref{sec:reflection}). Both
pairs \eqref{eq:shift} of an elliptic wave are counted; for a
rational-exponential wave the condition \(v\to0\) selects one of them.

\section{Pole structure and completeness}\label{sec:completeness}

This section proves Theorem~\ref{thm:complete} below: for odd \(p\),
and for even polynomials \(P\), every nonconstant meromorphic solution
of \eqref{eq:main} is rational, rational in an exponential, or
elliptic, and its poles form a single orbit under its periods.
The proof has three steps. We first compute the Laurent expansion at
a pole and show that for odd \(p\) it is unique. For even \(P\) it is
not: one coefficient is left free, and we show that the energy, a
first integral of \eqref{eq:main}, fixes it, so that Eremenko's
theorem applies on each energy level. Finally we show that any two
poles of one solution differ by a period.

Throughout, a meromorphic function is single-valued and meromorphic
on the whole complex plane. Let \(\Wclass\) consist of rational
functions of \(z\), rational functions of \(e^{\kappa z}\) with
\(\kappa\ne0\), and elliptic functions.
The period group of a nonconstant meromorphic function is its
\emph{full} group of periods: trivial for a rational function,
infinite cyclic for a rational function of \(e^{\kappa z}\), where
we choose \(\kappa\) so that \(2\pi i/\kappa\) generates it, and the
period lattice for an elliptic function.
For the general analytic background, see \citet{Eremenko1982}.

The theorem collects the analytic input to the enumeration.
Its odd-order case is contained in
\citet[Theorem~1]{YuanLiQi2014}.
For the even-order case, the first-integral mechanism is that of
\citet[Lemmas~3--4]{EremenkoLiaoNg2009}; here it is written uniformly
for every even polynomial \(P\). We give the proofs in full because
Sections~\ref{sec:systems} and~\ref{sec:counting} use the Laurent
recurrence and the energy dependence in explicit form.

\begin{theorem}[Completeness in the covered operator classes]
\label{thm:complete}
Assume either that \(p\) is odd or that \(P\) is an even polynomial.
Every nonconstant meromorphic solution of \eqref{eq:main} belongs to
\(\Wclass\). Every pole has order \(p\), and the pole set is one coset
of the full period group. Consequently there is one pole in the
rational case, one pole modulo the primitive period in the simply
periodic case, and one pole on the full elliptic period torus.

For fixed \(P\) of odd degree, there is at most one nonconstant solution
with its pole at zero. For even \(P\), the same statement holds after
fixing the conserved energy. The only entire solutions are
\(v=0\) and \(v=-2a_0\).
\end{theorem}

An even polynomial is one in which every odd power of \(\rho\) has
coefficient zero, so that \(P(D)\) contains only even derivatives.
The hypothesis does not cover an arbitrary operator of even degree,
and it cannot: Section~\ref{sec:open} gives a mixed second-order
operator with a meromorphic solution outside \(\Wclass\).

\subsection{Pole order and the Laurent recurrence}

Write a pole expansion at zero as
\begin{equation}\label{eq:laurent}
 v(z)=\sum_{n\ge0}c_nz^{n-p}.
\end{equation}
The pole order is forced by dominant balance. A pole of order \(r\)
makes \(P(D)v\) of order \(r+p\) and \(v^2\) of order \(2r\); the two
can cancel only if \(r=p\). Comparing the leading coefficients gives
\begin{equation}\label{eq:leadingpole}
 c_0=c_\ast:=-2(-1)^p R_p,\qquad
 R_p=\frac{(2p-1)!}{(p-1)!}.
\end{equation}
For the remaining coefficients, set
\begin{equation}\label{eq:indicial}
 I_p(n)=\prod_{j=p}^{2p-1}(n-j)+c_\ast.
\end{equation}
The product is the factor produced by \(D^p\) acting on \(z^{n-p}\),
and \(c_\ast\) is the factor produced by the cross term of
\(c_nz^{n-p}\) with the leading pole in \(v^2/2\).
With \(c_n=0\) for \(n<0\), comparison of the coefficients of
\(z^{n-2p}\) in \eqref{eq:main} gives
\begin{equation}\label{eq:recurrence}
 I_p(n)c_n=
 -\sum_{j=0}^{p-1}a_j
       \fall{(n-2p+j)}{j}\,c_{n-p+j}
 -\frac12\sum_{k=1}^{n-1}c_kc_{n-k},\qquad n\ge1.
\end{equation}
The underline denotes a falling factorial.
The right side contains only coefficients with index less than \(n\),
so \eqref{eq:recurrence} determines \(c_n\) whenever \(I_p(n)\ne0\).
The positive integer zeros of \(I_p\) are the \emph{resonances} of the
recurrence: at such an index the recurrence leaves \(c_n\)
undetermined if its right side vanishes, and has no solution
otherwise.

If \(p\) is odd, then \(c_\ast=2R_p\), and there is no positive
integer resonance.
For integers \(1\le n<p\), the product in \eqref{eq:indicial}
is negative with absolute value less than \(R_p\).
For \(p\le n\le2p-1\) it is zero, and for \(n\ge2p\) it is positive.
Hence \(I_p(n)\ne0\) for every positive integer \(n\), and there is
exactly one formal Laurent expansion at the marked pole.

To compare directly with \citet{YuanLiQi2014}, put \(w=-v/2\).
Then \(w^{(p)}+\sum_{j<p}a_jw^{(j)}=w^2\).
Their pole order is \(p\), their number of principal parts is one,
and each lower linear monomial has weighted pole order
\(p+j<2p\). Their degree and dominant-term hypotheses therefore hold.
We nevertheless apply Eremenko's theorem directly, since the same
theorem will be applied to an energy level in the even case.

Theorem~2 of \citet{Eremenko2006} states that an autonomous polynomial
differential equation with finitely many possible formal pole
expansions and a single monomial of highest total degree has all its
meromorphic solutions in \(\Wclass\).
In \eqref{eq:main} the only monomial of degree two is \(v^2/2\), and
for odd \(p\) we have just seen that there is exactly one formal pole
expansion. Both hypotheses hold at every odd order, for every
coefficient vector \(a\).

\subsection{The energy argument for an even operator}

Let \(p=2m\) and
\begin{equation}\label{eq:evenP}
 P(D)=D^{2m}+\sum_{j=0}^{m-1}b_jD^{2j}.
\end{equation}
Now \(c_\ast=-2R_p\), and there is exactly one positive integer
resonance.
Indeed the product in \eqref{eq:indicial} is at most \(R_p\)
for \(0\le n<p\), is zero for \(p\le n\le2p-1\), and is
strictly increasing for \(n\ge2p\); its value at \(3p\) is
\(2R_p=-c_\ast\). The only positive integer zero of \(I_p\) is
therefore \(n=3p\), and the possible free Laurent coefficient is
\[
 h=c_{3p}=[z^{2p}]v.
\]
If the recurrence is compatible at \(n=3p\), the equation has a
one-parameter family of formal expansions at the pole, one for each
value of \(h\), and Eremenko's theorem does not apply to
\eqref{eq:main} itself. The remedy is a first integral whose value
determines \(h\).

For \(j\ge1\), define the quadratic differential expression
\begin{equation}\label{eq:energyterms}
 \mathcal B_j(v)=
 \sum_{\ell=0}^{j-2}(-1)^\ell
       v^{(\ell+1)}v^{(2j-1-\ell)}
 +\frac{(-1)^{j-1}}2\bigl(v^{(j)}\bigr)^2,
\end{equation}
where the sum is empty when \(j=1\).
Termwise differentiation gives
\(D\mathcal B_j=v'v^{(2j)}\): every even derivative, multiplied by
\(v'\), is an exact derivative. Since \(P\) contains only even
derivatives, it follows that
\begin{equation}\label{eq:energy}
 \mathcal H_P(v)=
 \mathcal B_m(v)+\sum_{j=1}^{m-1}b_j\mathcal B_j(v)
 +\frac{b_0}{2}v^2+\frac16v^3
\end{equation}
satisfies
\begin{equation}\label{eq:energyidentity}
 D\mathcal H_P=v'E_P(v).
\end{equation}
In particular \(\mathcal H_P\), the energy, is constant on each
solution.
The case with no intermediate derivatives is the first integral
used by \citet{EremenkoLiaoNg2009}.
The familiar Kawahara first integral is the case \(m=2\)
\citep{DeminaKudryashovKawahara2010,YeEtAl2022};
the corresponding order-six construction appears in
\citet{Dang2023Even}.

The energy identity has two consequences for the Laurent expansion.
The first is that the resonance is always compatible.
To see this, construct \eqref{eq:laurent} through index \(3p-1\).
The first possible residual of \eqref{eq:main} is
\(\Omega_pz^p\).
Multiplication by \(v'\) gives residue \(-p c_\ast\Omega_p\).
By \eqref{eq:energyidentity} this is the residue of a derivative
of a Laurent series, and is therefore zero. Hence \(\Omega_p=0\),
and \(h\) may be chosen arbitrarily.

The second consequence is that the energy determines \(h\).
On the full formal solution the constant energy has the form
\begin{equation}\label{eq:energyh}
 \mathcal H_P=\Lambda_p h+\eta_p(P),\qquad
 \Lambda_p=-p c_\ast I_p'(3p)\ne0.
\end{equation}
To find the coefficient, linearize the highest-weight part of
\eqref{eq:energyidentity} about \(c_\ast z^{-p}\), with perturbation
\(z^{n-p}\).
If the corresponding energy term is \(L(n)z^{n-3p}\), then
\[
 (n-3p)L(n)=-p c_\ast I_p(n).
\]
Taking \(n\to3p\) gives the coefficient in \eqref{eq:energyh}.
Lower derivative terms have different powers of \(z\) and cannot
contribute to this coefficient. Moreover,
\[
 I_p'(3p)=2R_p\sum_{j=p+1}^{2p}\frac1j>0.
\]
Fixing the energy thus fixes \(h\), and all subsequent Laurent
coefficients follow from \eqref{eq:recurrence}.
This is the resonance-removal step in
\citet[Lemma~4]{EremenkoLiaoNg2009}, with an explicit coefficient
for the present operator family; see also
\citet[Section~3.2]{ConteNgWu2018} and
\citet[Section~5]{ConteMusetteNgWu2025}.

We can now apply Eremenko's theorem, not to \eqref{eq:main} but to
the level equation \(\mathcal H_P(v)=H\) with \(P\) and \(H\) fixed.
This equation is autonomous and polynomial, and its unique monomial of
highest degree is \(v^3/6\).
It has exactly one formal pole expansion: differentiating it and
dividing by the nonzero Laurent series \(v'\) returns \(E_P(v)=0\),
so every expansion belongs to the family above, and \eqref{eq:energyh}
selects the value of \(h\).
Every meromorphic solution of \eqref{eq:main} lies on one energy level
and therefore belongs to \(\Wclass\).

\subsection{The period group and the entire solutions}

It remains to show that the poles of one solution form a single orbit
under its periods.
For either covered case, take two poles \(z_1,z_2\) of the same
solution. The functions \(v(z+z_1)\) and \(v(z+z_2)\) satisfy the
same equation and, in the even case, have the same energy.
Their Laurent series at zero coincide, so the identity theorem gives
\[
 v(z+z_1)=v(z+z_2).
\]
Every difference of poles is a period. Conversely, every period
translates a pole to a pole. The period group is discrete: a sequence
of nonzero periods tending to zero would force \(v'\equiv0\).
The pole set is therefore one coset of that group.
The same argument proves the uniqueness statements of the theorem:
two nonconstant solutions with the same operator, the same energy in
the even case, and a pole at zero have the same Laurent series there.

Finally, an entire solution in \(\Wclass\) must be constant.
In the rational case it is a polynomial, and the
highest degree of its square cannot be cancelled in \eqref{eq:main}
unless the polynomial is constant.
In the rational-exponential case it is a Laurent polynomial in
\(t=e^{\kappa z}\); the operator \(P(\kappa t\partial_t)\) preserves
each exponent, whereas squaring doubles an extreme nonzero exponent.
An entire elliptic function is constant by compactness.
The constants solve \(a_0v+v^2/2=0\), which gives \(v=0\) and
\(v=-2a_0\).
This completes the proof of Theorem~\ref{thm:complete}.

The finite systems of the next section are derived from the one-pole
structure alone. The analytic facts of this section are used only to
show that, in the covered cases, those systems miss no nonconstant
meromorphic solution. The systems themselves never require the convergence of a
formal series.

\section{Three finite algebraic systems}\label{sec:systems}

Theorem~\ref{thm:complete} says that, in the covered cases, a
nonconstant meromorphic solution is one of three kinds and has one
pole modulo its periods. Each kind is a finite-dimensional family of
functions: a polynomial in \(1/z\), a polynomial in a fixed simply
periodic function with one pole, or a constant plus a linear
combination of \(\wp\) and its derivatives. Substituting the general member of each family
into \eqref{eq:main} therefore gives finitely many polynomial
equations in its coefficients. This section writes down the three
systems, fixes their normalizations, and proves that they are exactly
equivalent to the one-pole solutions (Theorem~\ref{thm:systems}).
They are instances of the principal-part construction of
\citet{DeminaKudryashov2011,DeminaKudryashovElliptic2011}.

The systems all have the same shape. Let \(V\) be a polynomial
expression in some variables, and let \(\delta\) be a differentiation
rule for those variables, so that the polynomials in them form a
differential algebra with derivation \(\delta\). Write
\[
 \mathcal E_P(V;\delta)=
 \delta^pV+\sum_{j=0}^{p-1}a_j\delta^jV+\frac12V^2.
\]
Each system below is the polynomial identity
\(\mathcal E_P(V;\delta)=0\) for an appropriate \(V\) and \(\delta\).
A polynomial identity means the vanishing of every coefficient.

\subsection{Rational solutions}

A rational solution whose only pole is at zero, of order \(p\), and
which has no nonconstant polynomial part at infinity is a polynomial
of degree \(p\) in \(T=1/z\). Since \(DT=-T^2\), the equation becomes
\begin{equation}\label{eq:rational-system}
 V_R(T)=b+\sum_{k=1}^{p}r_kT^k,\qquad r_p=c_\ast,\qquad
 \mathcal E_P(V_R;-T^2\partial_T)=0.
\end{equation}
The left side of the identity is a polynomial of degree at most
\(2p\) in \(T\), so this is a finite system in
\(b,r_1,\ldots,r_{p-1}\).
No rational solution has a nonconstant polynomial part at infinity:
such a part would have its degree doubled by \(v^2\), while the
linear differential expression cannot reach that degree.
Consequently \eqref{eq:rational-system} is the full rational
one-pole system.

\subsection{Rational functions of an exponential}

For \(q=\kappa^2\ne0\), define
\begin{equation}\label{eq:T}
 T_q(z)=\frac{\kappa}{2}\coth\frac{\kappa z}{2},
 \qquad
 DT_q=\frac q4-T_q^2.
\end{equation}
The expression is independent of the choice of sign of \(\kappa\).
It has period \(2\pi i/\kappa\) and one pole per period, at zero,
with residue one; it plays the role that \(1/z\) plays in the
rational case.
The finite system is
\begin{equation}\label{eq:exponential-system}
 V_S(T)=\sum_{k=0}^{p}u_kT^k,\qquad
 u_p=c_\ast,\qquad q\ne0,\qquad
 \mathcal E_P\left(V_S;\left(\frac q4-T^2\right)\partial_T\right)=0.
\end{equation}

To see that every one-pole solution rational in an exponential has
this form, write \(v=R(t)\), \(t=e^{\kappa z}\).
The rational function \(R\) has finite limits at both \(t=0\) and
\(t=\infty\). If it had a pole at either end, an extreme Laurent
monomial would be doubled by \(R^2\) and could not cancel against
\(P(\kappa t\partial_t)R\).
After a translation, its only finite pole is at \(t=1\), of order \(p\).
Such a function is a degree-\(p\) polynomial in \(T_q\).
The logistic coordinate
\[
 Q=\frac1{1+e^{-z}},\qquad DQ=Q(1-Q),
\]
is an equivalent coordinate after fixing the rate and translating the
pole; \(Q(z+i\pi)=T_1(z)+1/2\).
Polynomial logistic representations and their use in constructing
travelling waves are developed by
\citet{Kudryashov2012,Kudryashov2013,Kudryashov2015}.

The rational system is the \(q=0\) specialization of
\eqref{eq:exponential-system}, with \(T_0=1/z\).
We keep \(q=0\) and \(q\ne0\) distinct so that this degeneration is
never counted twice.

\subsection{Elliptic solutions}

Use the Weierstrass algebra
\begin{equation}\label{eq:weierstrass}
 Y^2=4X^3-g_2X-g_3,\qquad
 \mathscr D X=Y,\qquad
 \mathscr D Y=6X^2-\frac{g_2}{2},\qquad
 \Delta=g_2^3-27g_3^2\ne0.
\end{equation}
Its analytic realization is \(X=\wp(z;g_2,g_3)\),
\(Y=\wp'(z;g_2,g_3)\); see \citet[Chapter~23]{DLMF}.
Set
\begin{equation}\label{eq:elliptic-system}
 V_E=b+\sum_{j=0}^{p-2}t_j\mathscr D^jX,\qquad
 t_{p-2}=-s_p,\qquad
 s_p=\frac{2(2p-1)!}{((p-1)!)^2},
 \qquad
 \mathcal E_P(V_E;\mathscr D)=0.
\end{equation}
Using \eqref{eq:weierstrass}, every polynomial in \(X\) and \(Y\)
reduces to the normal form \(F(X)+YG(X)\). The system consists of
the coefficients of \(F\) and \(G\) set equal to zero; all of them
are polynomials with rational coefficients in the displayed unknowns.

The representation is the classical one. An elliptic function with
one pole has zero residue there.
The functions \(1,\wp,\wp',\ldots,\wp^{(p-2)}\) have distinct
principal parts and span the space of elliptic functions whose only
pole is at zero, of order at most \(p\).
Subtracting the indicated principal part from a one-pole solution
leaves an entire elliptic function, hence a constant.
The leading coefficient in \eqref{eq:elliptic-system} follows from
\[
 D^{p-2}\wp=(-1)^{p-2}(p-1)!z^{-p}+O(1).
\]
The nonsingularity condition \(\Delta\ne0\) is retained throughout;
the singular cubics correspond to the rational and simply periodic
systems, which we have already written down.

\begin{theorem}[Finite matching and exhaustion]\label{thm:systems}
For every \(p\ge2\), the systems
\eqref{eq:rational-system}, \eqref{eq:exponential-system}, and
\eqref{eq:elliptic-system} are necessary and sufficient for the
respective one-pole solutions of \eqref{eq:main}, with the pole at
zero and with the stated normalizations.
For the operator classes of Theorem~\ref{thm:complete}, these systems,
together with the constant solutions, exhaust all globally meromorphic
solutions.

The systems describe exceptional parameter values and continuous
families as well as isolated solutions. Their sufficiency follows
from substitution, without any genericity assumption.
\end{theorem}

\begin{proof}
Necessity is the principal-part argument given for each kind above.
Conversely, a solution of a coefficient system gives a globally
meromorphic function by its displayed realization, and the chain rule
gives \(E_P(v)=0\).
The leading coefficient \(c_\ast\ne0\) ensures a nonconstant function.
In the elliptic case every nonsingular pair \((g_2,g_3)\) has a
Weierstrass realization, and the normal form \(F(X)+YG(X)\) is unique.
The exhaustion statement is Theorem~\ref{thm:complete}.
\end{proof}

\subsection{Eliminating the profile coefficients}

The systems above treat the operator and the profile coefficients as
unknowns together, which is the form the enumeration in
Section~\ref{sec:counting} requires. For a fixed operator they can be
reduced further, because the Laurent recurrence already determines
the principal part of any solution. Compute
\[
 r_k=c_{p-k}\quad(1\le k\le p),\qquad s=c_p
\]
from \eqref{eq:recurrence}.
These coefficients occur before the even-order resonance and are
determined by \(P\) at either parity.
In the simply periodic algebra define
\[
 S_k(T;q)=\frac{(-1)^{k-1}}{(k-1)!}
 \left[\left(\frac q4-T^2\right)\partial_T\right]^{k-1}T.
\]
It has principal part \(z^{-k}\).
Its regular Laurent constant is
\[
 d_k(q)=
 \begin{cases}
 0,&k\text{ odd},\\
 -B_kq^{k/2}/k!,&k\text{ even},
 \end{cases}
\]
where \(B_k\) denotes a Bernoulli number.
Thus the unique candidate with the required principal part and
constant is
\begin{equation}\label{eq:genuszero-candidate}
 V_0(T;P,q)=s+\sum_{k=1}^{p}r_k\bigl(S_k(T;q)-d_k(q)\bigr).
\end{equation}
Substituting it into \eqref{eq:main} leaves polynomial equations in
the single unknown \(q\).

For the elliptic candidate, first require \(r_1=0\), since an
elliptic function with one pole has no residue there, and put
\begin{equation}\label{eq:elliptic-candidate}
 \widetilde V_E=
 s+\sum_{k=2}^{p}r_k
 \left(\frac{(-1)^{k-2}}{(k-1)!}\mathscr D^{\,k-2}X
       -e_k(g_2,g_3)\right),
\end{equation}
where \(e_k\) is the regular Laurent constant of the corresponding
derivative of \(\wp\).
It is a polynomial in \(g_2,g_3\), computed from
\eqref{eq:weierstrass}.
The remaining equations involve only \(g_2,g_3\), with
\(\Delta\ne0\).
This is the coefficient-comparison construction used in
\citet{DeminaKudryashov2011,DeminaKudryashov2014}.

For even \(P\) the equation \(E_P(v)=0\) does not see the energy: a
candidate may solve it on any energy level. To classify a fixed
level, compute \(h\) from \eqref{eq:energyh} for the prescribed
energy and append
\begin{equation}\label{eq:selecth}
 [z^{2p}]V_0(T_q;P,q)=h
 \quad\text{or}\quad
 [z^{2p}]\widetilde V_E(\wp,\wp';P,g_2,g_3)=h.
\end{equation}
Both are polynomial equations obtained from finite Laurent expansions.
With this equation included, the reduced systems classify a fixed
energy level: at fixed \(P\) of odd degree, or fixed \((P,h)\) in the
even case, there is at most one nonconstant solution with its pole at
zero.

When all prescribed data are exact algebraic numbers, that is, the
coefficients of \(P\) and, in the even case, the energy, these are
finitely many polynomial equations in one or two unknowns.
Polynomial gcds in the one-variable case and elimination in the
multivariable case decide whether a solution exists.
A nonvanishing condition such as \(\Delta\ne0\) is imposed by
adjoining \(u\Delta-1=0\). Elimination with case distinctions
describes instead the operators, or the marked data \((P,h)\), that
admit each kind of solution, as a finite union of systems of
polynomial equations and nonvanishing conditions
\citep{CoxLittleOShea2005}.

\section{The count of rational-exponential pairs}\label{sec:counting}

We count the pairs \((P,v)\) of Section~\ref{sec:setting} for all monic
\(P\) of degree \(p\). Whether an equation can have other meromorphic
waves is a separate question (Section~\ref{sec:completeness}); nothing
in this section or the next depends on the answer.

By Section~\ref{sec:systems}, a normalized profile is a polynomial of
degree \(p\) in \(Q=1/(1+e^{-z})\), without constant term because
\(v\to0\) as \(e^z\to0\). Since \(D^jQ=(-1)^j j!Q^{j+1}\) plus lower
powers of \(Q\), the functions \(Q,DQ,\ldots,D^{p-1}Q\) span the
degree-\(p\) polynomials in \(Q\) with zero constant term, and the
leading pole coefficient \eqref{eq:leadingpole} gives the unique form
\begin{equation}\label{eq:normalized-A}
 v_A=s_pA(D)Q,\qquad
 A(\rho)=\rho^d+\alpha_1\rho^{d-1}+\cdots+\alpha_d,\qquad d=p-1,
\end{equation}
with the scale \(s_p\) of \eqref{eq:leadingpole}; the unknowns of the
count are \(\alpha_1,\ldots,\alpha_d\). This is the standard
singular-part operator form: \(Q=1/2+D\log\cosh(z/2)\), so \(v_A\) is a
constant plus a linear differential operator applied to
\(\log\cosh(z/2)\); such forms, including the KS case, are discussed by
\citet{ConteMusette2003}.

\subsection{Discrete convolution and the matching equations}

For polynomials \(F,G\), let \(F\star G\) be the polynomial with
\((F\star G)(n)=\sum_{j=1}^{n-1}F(j)G(n-j)\) for all integers \(n\ge1\).
It exists because, by the power-sum formulas, the right side is a
polynomial in \(n\). Write \(S_A=A\star A\). For monic \(A\) of degree
\(d\),
\begin{equation}\label{eq:Sleading}
 \deg S_A=2d+1,\qquad
 [\rho^{2d+1}]S_A=\frac{(d!)^2}{(2d+1)!}=\frac{2}{s_p};
\end{equation}
the leading coefficient is the beta integral \(\int_0^1t^d(1-t)^d\dd t\).

\begin{proposition}[Polynomial divisibility]\label{prop:divisibility}
The profile \eqref{eq:normalized-A} solves \eqref{eq:main} for some
monic \(P\) if and only if \(A\) divides \(S_A\).
The operator is unique and is
\begin{equation}\label{eq:operator-A}
 P_A=\frac{s_pS_A}{2A}.
\end{equation}
\end{proposition}

\begin{proof}
In the half-plane \(\operatorname{Re}z<0\), the logistic function is a
geometric series in \(e^z\), and applying \(A(D)\) multiplies the
\(n\)th term by \(A(n)\): \(A(D)Q=\sum_{n\ge1}(-1)^{n-1}A(n)e^{nz}\).
Squaring gives the rational-function identity
\begin{equation}\label{eq:discrete-square}
 (A(D)Q)^2=-S_A(D)Q.
\end{equation}
The map \(F\mapsto F(D)Q\) is injective: its vanishing would give
\(F(n)=0\) for every positive integer \(n\). Consequently
\eqref{eq:main} is equivalent to \(2P_AA=s_pS_A\), which has a
polynomial solution \(P_A\) exactly when \(A\) divides \(S_A\). Equation
\eqref{eq:Sleading} shows that \(P_A\) is monic of degree \(p\).
\end{proof}

Divide \(S_A\) by the monic polynomial \(A\):
\begin{equation}\label{eq:remainder}
 S_A=AC_A+R_A,\qquad
 R_A=\sum_{r=0}^{d-1}R_r(\alpha_1,\ldots,\alpha_d)\rho^r.
\end{equation}
The matching equations are \(R_0=\cdots=R_{d-1}=0\): \(d\) polynomial
equations in the \(d\) unknowns \(\alpha_1,\ldots,\alpha_d\). Their
solutions are in bijection with the normalized pairs \((P_A,v_A)\); we
write \(\Xp\) for the set of solutions. For \(p=2\) and \(A=\rho+a\)
there is the single equation \(R_0=-a(a-1)(a+1)/6=0\) of the
Introduction. Multiplicities in \(\Xp\) refer to these equations.

\paragraph{Simple solutions and multiplicity.}
A solution is \emph{simple} if the Jacobian matrix
\((\partial R_r/\partial\alpha_j)\) is invertible there, like a simple
root, \(f(a)=0\), \(f'(a)\ne0\). Every isolated solution, that is, one
that is not a limit of other solutions, has a \emph{multiplicity}: the
number of solutions near it of the system \(R_r=\varepsilon_r\) for
generic small \(\varepsilon\in\CC^d\), as the double root of \(x^2=0\)
becomes the two roots of \(x^2=\varepsilon\). It is one exactly when the
solution is simple, and it equals the dimension over \(\CC\) of the
convergent power series at that point modulo the combinations
\(\sum_rh_rR_r\) with power series \(h_r\). For finitely many solutions,
the sum of their multiplicities, the \emph{count with multiplicity}, is
the dimension of the polynomials modulo the polynomial combinations
\citep[Sections~4.3 and~5.2]{ArnoldGuseinZadeVarchenko1985},
\citep[Chapter~4, \S2]{CoxLittleOShea2005}; it is the number of distinct
pairs exactly when every solution is simple (see Section~\ref{sec:ks}).

\paragraph{How many solutions can there be?}
Before counting, we ask how many solutions such a system can have at
most. The answer is already \(N_p\). Give \(\alpha_j\) weight \(j\) and
\(\rho\) weight one, so that \(A\) is weighted-homogeneous of weight
\(d\): replacing the roots of \(A\) by \(\lambda\) times themselves
multiplies \(\alpha_j\) by \(\lambda^j\). The polynomial \(S_A\) has
weight at most \(2d+1\), and its part of weight exactly \(2d+1\) is the
\emph{continuous} convolution
\begin{equation}\label{eq:continuous}
 H_A(\rho)=\int_0^\rho A(t)A(\rho-t)\dd t,\qquad
 \int_0^\rho t^i(\rho-t)^j\dd t=\frac{i!\,j!}{(i+j+1)!}\,\rho^{i+j+1}.
\end{equation}
Indeed these are the highest weighted terms of the power sums that
define \(S_A\). Monic division preserves the highest weighted part:
\(A\) is weighted-homogeneous and monic, so each step of long division
subtracts a multiple of \(A\) of the same weight as the term being
removed. Consequently the matching equation \(R_r\) has weight at most
\(2d+1-r\), and
\begin{equation}\label{eq:weights}
 \text{the part of weight }2d+1-r\text{ of }R_r
 \text{ is the coefficient of }\rho^r\text{ in }\remA H_A .
\end{equation}
In words, the discrete convolution \(S_A\) is a lower-order perturbation
of the continuous convolution \(H_A\), and far from the origin the
matching equations look like the continuous problem \(A\mid H_A\).
Solutions of the matching equations could escape to infinity only along
solutions of that problem. \emph{Rigidity} \((\mathrm R_d)\) is the
statement that it has none except the trivial one:
\begin{equation}\label{eq:rigidity}
 A\text{ monic},\quad\deg A=d,\quad A\mid H_A
 \quad\Longrightarrow\quad A(\rho)=\rho^d.
\end{equation}

\begin{proposition}[Weighted B\'ezout bound]\label{prop:bezout}
Let \(p\ge2\) and \(d=p-1\). The matching equations are \(d\) equations
of weighted degrees at most \(2d+1,2d,\ldots,d+2\) in \(d\) unknowns of
weights \(1,2,\ldots,d\). Their isolated solutions, counted with
multiplicity, number at most
\[
 \frac{(d+2)(d+3)\cdots(2d+1)}{d!}=\binom{2d+1}{d}
 =\binom{2p-1}{p-1}=N_p .
\]
If rigidity \((\mathrm R_d)\) holds, that is, if the parts of weight
\(2d+1-r\) of the \(R_r\), which by \eqref{eq:weights} are the
coefficients of \(\remA H_A\), have no common zero other than the
origin, then the matching equations have finitely many solutions, and
counted with multiplicity there are exactly \(N_p\) of them. They are
then \(N_p\) distinct solutions exactly when all of them are simple.
\end{proposition}

For \(p=2,3\) the bound is \(3/1=3\) and \((4\cdot5)/(1\cdot2)=10\). The
substitution \(\alpha_j=y_j^j\) turns weights into ordinary degrees and
counts every solution \(d!\) times; the proof, from B\'ezout's inequality
for isolated solutions \citep[Chapter~VIII]{Mondal2021} and B\'ezout's
theorem \citep{CoxLittleOShea2005}, is Theorem~2.1 of the companion
note. The proof of Theorem~\ref{thm:prime} does not use the proposition:
there the labels give the solutions, and rigidity is a by-product. The
proposition is what remains at other orders (Section~\ref{sec:further}).

For waves, rigidity is a statement about rational solutions. The
dictionary is the Laplace transform, restricted to polynomials.

\begin{lemma}[Laplace correspondence]\label{lem:laplace}
Let \(F\mapsto F^\sharp\) be the linear map on polynomials with
\((\rho^k)^\sharp=k!\,z^{-k-1}\). Then \(D(F^\sharp)=(-\rho F)^\sharp\),
and \(F^\sharp G^\sharp=(F\ast G)^\sharp\) for the continuous convolution
\((F\ast G)(\rho)=\int_0^\rho F(t)G(\rho-t)\dd t\); in particular
\((A^\sharp)^2=(H_A)^\sharp\).
Let \(P\) be monic of degree \(p\), and let \(V\) be a nonconstant
rational solution of \eqref{eq:main} whose only pole is at \(z=0\).
Then \(V\) vanishes at infinity, \(P(0)=0\), and
\[
 V=(-1)^{p-1}s_pA^\sharp,\qquad A\mid H_A,\qquad
 P(\rho)=(-1)^p\,\frac{s_p}{2}\,\frac{H_A}{A}(-\rho),
\]
with \(A\) monic of degree \(p-1\). Conversely, every monic \(A\) of
degree \(p-1\) that divides \(H_A\) gives such a solution. The residue
of \(V\) at zero is \((-1)^{p-1}s_pA(0)\).
\end{lemma}

\begin{proof}
Both rules are checked on monomials; the second is the integral in
\eqref{eq:continuous}. By \eqref{eq:rational-system}, \(V=b+f\) with a
constant \(b\) and with \(f\to0\) at infinity. The constant part of
\eqref{eq:main} is \(b\,(P(0)+b/2)=0\), and the slowest decaying term of
\(f\) enters the rest of the equation with the factor \(P(0)+b\), which
must vanish. Hence \(b=P(0)=0\). By \eqref{eq:leadingpole} the pole has
order \(p\) and leading coefficient \(c_\ast=(-1)^{p-1}(p-1)!\,s_p\), so
\(V=cA^\sharp\) with \(c=(-1)^{p-1}s_p\) and \(A\) monic of degree
\(p-1\). By the two rules, \eqref{eq:main} for \(V\) reads
\(\bigl(cP(-\rho)A+c^2H_A/2\bigr)^\sharp=0\), and the map \(\sharp\) is
injective. This gives the divisibility, the formula for \(P\), and the
converse: the \(P\) defined by the formula is monic of degree \(p\),
because \(H_A\) has the leading coefficient \eqref{eq:Sleading} of
\(S_A\). The residue is the coefficient \(cA(0)\) of \(z^{-1}\).
\end{proof}

For \(A=\rho^{p-1}\) the lemma gives \(V=c_\ast z^{-p}\) and \(P=D^p\).
Rigidity \((\mathrm R_{p-1})\) is therefore the statement that, for
every monic \(P\) of degree \(p\), equation \eqref{eq:main} has no other
nonconstant rational solution with a single pole.

\subsection{The prime-index theorem}\label{sec:prime}

\begin{theorem}[Prime-index enumeration]\label{thm:prime}
Let \(\ell=2p-1\) be prime, and put \(d=p-1\).
\begin{enumerate}
\item[\textup{(a)}] The matching equations have exactly
\(\binom{\ell}{d}=N_p\) solutions over \(\CC\), and every solution is
simple. Thus there are exactly \(N_p\) rational-exponential pairs, and
they are distinct.
\item[\textup{(b)}] The coefficients of every solution are algebraic
numbers, and they are \(\ell\)-integral: they have no \(\ell\) in a
denominator and can be reduced modulo \(\ell\), in the way in which
\(1/11\) becomes \(1\) and \(i\) becomes \(2\) or \(3\) modulo \(5\).
Once a consistent way of reducing algebraic numbers is fixed
(Section~\ref{sec:prime-proof}), reduction modulo \(\ell\) is a
bijection between the solutions and the \(d\)-element subsets \(L\) of
the field \(\FF_\ell\) of integers modulo \(\ell\):
\begin{equation}\label{eq:prime-labels}
 \overline A(\rho)=\prod_{s\in L}(\rho-s).
\end{equation}
\item[\textup{(c)}] Rigidity \((\mathrm R_d)\) holds.
\end{enumerate}
\end{theorem}

We call \(L\) the \emph{label} of the pair. It depends on the way of
reducing algebraic numbers, an embedding fixed in the proof, only as
\(i\) is \(2\) or \(3\) modulo \(5\): another embedding gives a solution
the label that the first one gives to a conjugate solution. Whether a
label contains \(0\), and the labels of solutions with rational
coefficients, do not depend on it. The labels describe reductions of
algebraic coefficients; they do not assert that the complex roots of
\(A\), or the characteristic roots of \(P_A\), are integers.

\paragraph{Example: order three.}\label{sec:ks-rational}
Let \(p=3\), \(s_3=60\), and \(A(\rho)=\rho^2+a\rho+b\), so that
\((a,b)=(\alpha_1,\alpha_2)\). The profile and monic operator are
\begin{align}
 \frac{v_A}{60}&=2Q^3-(a+3)Q^2+(a+b+1)Q, \label{eq:ksprofile}\\
 P_A(\rho)&=\rho^3+4a\rho^2+(a^2+19b)\rho-a^3+7ab-5a-30b.
   \label{eq:ksoperator}
\end{align}
The remainder of \(S_A\) is \((\widehat R_1\rho+\widehat R_0)/30\) with
\begin{equation}\label{eq:ksmatching}
 \widehat R_0=30R_0=ab(a^2-7b+5),\qquad
 \widehat R_1=30R_1=a^4-8a^2b+11b^2+10b-1.
\end{equation}
We solve \(\widehat R_0=\widehat R_1=0\) by cases on the factors of
\(\widehat R_0\). If \(b=0\), then \(a^4=1\). If \(a=0\), then
\((b+1)(11b-1)=0\). In the remaining case \(a^2=7b-5\), and the second
equation becomes \(4(b-2)(b-3)=0\). This exhausts the system and gives
Table~\ref{tab:ks}. The determinant
\(\det\partial(\widehat R_0,\widehat R_1)/\partial(a,b)\) is nonzero at
every point, so all \(N_3=10\) pairs are simple, as
Theorem~\ref{thm:prime}(a) states for \(2p-1=5\).

\begin{table}[htbp]
\centering
\caption{All ten rational-exponential pairs at order three. The label
is the set of roots of \(A\) modulo \(5\); for the two complex pairs it
depends on the choice of \(i\) modulo \(5\). For the equation of the
Introduction, \((b_1,b_2,b_3)=(1,\sigma,1)\), the coefficients \(B=4a\)
of \(\rho^2\) and \(\mu=a^2+19b\) of \(\rho\) in \eqref{eq:ksoperator}
are \(\sigma/\kappa\) and \(1/\kappa^2\) by \eqref{eq:coefficients}, so
\(\sigma^2=B^2/\mu=16a^2/(a^2+19b)\). In the rows \((\pm i,0)\) and
\((0,-1)\), \(\mu<0\) and the rate \(\kappa\) is imaginary. The last
column describes \(v_A\) on the real axis, at unit rate.}
\label{tab:ks}
\begin{tabular}{@{}cclccl@{}}
\toprule
\(a\)&\(b\)&Labels&\(|\sigma|\)&Determinant&Real profiles\\
\midrule
\(\pm1\)&\(0\)&\(\{0,4\}\), \(\{0,1\}\)&\(4\)&\(-24\)&Pulses\\
\(\pm i\)&\(0\)&\(\{0,2\}\), \(\{0,3\}\)&\(4\)&\(-16\)&Not real\\
\(0\)&\(-1\)&\(\{1,4\}\)&\(0\)&\(144\)&Front\\
\(0\)&\(1/11\)&\(\{2,3\}\)&\(0\)&\(576/121\)&Front\\
\(\pm3\)&\(2\)&\(\{3,4\}\), \(\{1,2\}\)&\(12/\sqrt{47}\)&\(-144\)&Fronts\\
\(\pm4\)&\(3\)&\(\{2,4\}\), \(\{1,3\}\)&\(16/\sqrt{73}\)&\(384\)&Fronts\\
\bottomrule
\end{tabular}
\end{table}

The ten labels are the ten two-element subsets of \(\{0,1,2,3,4\}\). By
\eqref{eq:ksprofile}, \(v_A/60\to b\) as \(z\to+\infty\), so a pair is a
pulse exactly when \(b=0\). The map \((a,b)\mapsto(-a,b)\) is the mirror
image of the Introduction, and it replaces a label \(L\) by \(-L\)
(Section~\ref{sec:reflection}). Eight of the ten pairs have real
coefficients, and all eight are smooth and bounded on the real axis: two
pulses and six fronts. For example, \(A=(\rho-1)(\rho-2)\) gives
\(v_A=120Q^3\) and \(P_A=(\rho-3)(\rho-4)(\rho-5)\). These are the
classical solutions in hyperbolic functions
\citep{KuramotoTsuzuki1976,Kudryashov1988,Kudryashov1989,Kudryashov1990,Kudryashov1991},
tabulated completely by \citet[Table~1]{KudryashovZargaryan1996} and
\citet[Table~1]{ConteMusette2009}. Their tables contain six cases
because they identify the two signs of the dispersion coefficient: the
four mirror pairs \((\pm a,b)\) and the two pairs with \(a=0\).
Exhaustion of this list follows from \citet{Eremenko2006};
\citet{Demina2020} classifies all irreducible algebraic invariants of
the same profile equation, and later rederivations by ansatz methods
coincide with these waves \citep{KudryashovSoukharev2009}. The pairs
\((\pm1,0)\) and \((\pm i,0)\) belong to the same equations, those with
\(\sigma^2=16\), at the rates \(\kappa=\pm1\) and \(\kappa=\pm i\)
(Section~\ref{sec:fixed}). For a real equation the second wave is thus
periodic in \(x\), and its real form has poles on the real axis.

\subsection{Proof of the prime-index theorem}\label{sec:prime-proof}

The idea was described in the Introduction. A reader who accepts the
theorem can go to Section~\ref{sec:labels}: the rest of the paper uses
only its statement, and the proof of Proposition~\ref{cor:endpoints} uses
\eqref{eq:prime-collapse} and fact \textup{(V)} below once more.

\paragraph{Reduction modulo \(\ell\) of algebraic numbers.}
We use four standard facts. Rational numbers without \(\ell\) in the
denominator form a ring \(\ZZ_{(\ell)}\). Think of \(\ell\) as a small
parameter \(\varepsilon\): the valuation is the order of a quantity in
\(\varepsilon\), and reduction modulo \(\ell\) sets \(\varepsilon=0\).

\emph{\textup{(V)} Valuation and reduction.} Let \(K\) be an algebraic
closure of the field \(\QQ_\ell\) of \(\ell\)-adic numbers. It is
algebraically closed, contains \(\QQ\), and carries the \(\ell\)-adic
valuation \(\operatorname{ord}\colon K\to\QQ\cup\{\infty\}\)
\citep[Chapter~III]{Koblitz1984},
\citep[Chapter~II, Theorem~(4.8)]{Neukirch1999}: the exponent of
\(\ell\) in a nonzero rational number, infinite only at \(0\), with
\begin{equation}\label{eq:valuation-rules}
 \operatorname{ord}(xy)=\operatorname{ord}(x)+\operatorname{ord}(y),\qquad
 \operatorname{ord}(x+y)\ge\min\{\operatorname{ord}(x),\operatorname{ord}(y)\}.
\end{equation}
We fix an embedding of the field \(\overline\QQ\subset\CC\) of algebraic
numbers in \(K\), unique up to an automorphism of \(\overline\QQ\), and
regard algebraic numbers as elements of \(K\). Those \(x\in K\) with
\(\operatorname{ord}(x)\ge0\) are \emph{\(\ell\)-integral}. They form a
ring \(\mathcal O\supset\ZZ_{(\ell)}\), and \emph{reduction modulo
\(\ell\)}, \(x\mapsto\overline x\), maps \(\mathcal O\) onto its residue
field \(\Bbbk\supset\FF_\ell\): it respects sums and products, it is the
usual reduction on \(\ZZ_{(\ell)}\), and \(\overline x=0\) exactly when
\(\operatorname{ord}(x)>0\). Polynomials are reduced coefficient by
coefficient. Division by a monic polynomial uses only sums and products
of coefficients: it commutes with substituting values for the
\(\alpha_j\) and with reduction, and it stays within
\(\ZZ_{(\ell)}[\alpha_1,\ldots,\alpha_d]\) or \(\mathcal O\). A monic
polynomial with \(\ell\)-integral roots has \(\ell\)-integral
coefficients.

\emph{\textup{(H)} Hensel's lemma.} Let \(G=(G_1,\ldots,G_d)\) be
polynomials in \(d\) unknowns with coefficients in \(\ZZ_{(\ell)}\). If
\(\overline a\in\FF_\ell^d\) solves the reduced system and its Jacobian
determinant is not zero there, then exactly one \(a\in\mathcal O^d\) has
\(G(a)=0\) and reduction \(\overline a\). This is Newton's method from
an integer vector that reduces to \(\overline a\), just as a simple root
at \(\varepsilon=0\) continues uniquely to a power series in
\(\varepsilon\): the multivariable Hensel lemma of \citet{ConradHensel},
applied in the finite extensions of \(\QQ_\ell\) in \(K\), which are
complete \citep[Chapter~II, Theorem~(4.8)]{Neukirch1999}.

\emph{\textup{(N)} The weak Nullstellensatz.} Polynomial equations with
rational coefficients that have a solution in \(\CC^d\) have one in
\(\overline\QQ^d\) \citep[Chapter~4, \S1]{CoxLittleOSheaIVA}.

\emph{\textup{(T)} Finite systems.} If \(d\) polynomial equations in
\(d\) unknowns with rational coefficients have only finitely many
solutions in \(K^d\), then these lie in \(\overline\QQ^d\), and they are
also the solutions in \(\CC^d\), with the Jacobian determinant nonzero
at the same ones. Indeed, by finiteness a nonzero polynomial in each
\(x_i\) alone is a combination of the equations
\citep[Chapter~5, \S3]{CoxLittleOSheaIVA}; finding one of bounded degree
is a linear problem over \(\QQ\), so one with rational coefficients
exists, and it confines \(x_i\) to its roots, in \(K\) as in \(\CC\).

\begin{proof}[Proof of Theorem~\ref{thm:prime}]
The proof has four steps: a congruence for \(S_A\) modulo \(\ell\); the
solution of the reduced equations; integrality, a dominant balance
showing that every solution can be reduced, and then reduces to one of
them; and lifting, Newton's method in the form of Hensel's lemma,
showing that each of them comes from exactly one solution. Solutions are
first taken in \(K^d\), with \(K\) as in \textup{(V)}; the last step
returns to \(\CC\).

\emph{Step 1: the congruence.} Induction using \((D-n)Q^n=-nQ^{n+1}\)
gives \((n-1)!\,Q^n=(-1)^{n-1}\prod_{j=1}^{n-1}(D-j)Q\). Write
\((A(D)Q)^2=\sum_{n=2}^{2p}\gamma_nQ^n\). Each \(\gamma_n\) is an
integer polynomial in the coefficients of \(A\), and
\(\gamma_{2p}=(d!)^2\). By \eqref{eq:discrete-square},
\[
 S_A(\rho)=\sum_{n=2}^{2p}\frac{(-1)^n\gamma_n}{(n-1)!}
       \prod_{j=1}^{n-1}(\rho-j).
\]
Only the last denominator, \((2p-1)!=\ell!\), contains \(\ell\).
Consequently \(\ell S_A\) has coefficients in
\(\ZZ_{(\ell)}[\alpha_1,\ldots,\alpha_d]\), and
\begin{equation}\label{eq:prime-collapse}
 \ell S_A(\rho)\equiv\frac{(d!)^2}{(\ell-1)!}\prod_{j=1}^{\ell}(\rho-j)
 \equiv(-1)^d(\rho^\ell-\rho)\pmod\ell.
\end{equation}
The scalar follows by pairing \(j\) with \(\ell-j\) in \((\ell-1)!\).
The product is \(\rho^\ell-\rho\) because, by Fermat's little theorem,
every residue \(s\) modulo \(\ell\) satisfies \(s^\ell=s\). The right
side does not depend on \(A\): this is what makes the prime index
special. Taking the parts of weight \(\ell=2d+1\) on both sides gives
the corresponding congruence for the continuous convolution,
\begin{equation}\label{eq:prime-continuous}
 \ell H_A(\rho)\equiv(-1)^d\rho^\ell\pmod\ell .
\end{equation}
Both congruences are between polynomials in \(\rho\) and the
\(\alpha_j\) with coefficients in \(\ZZ_{(\ell)}\), so they may be
evaluated at any \(\ell\)-integral values of the \(\alpha_j\) and of
\(\rho\), and then reduced.

\emph{Step 2: the reduced equations.} By \textup{(V)} the scaled
remainders \(\ell R_r\) have coefficients in \(\ZZ_{(\ell)}\), and by
\eqref{eq:prime-collapse} their reductions are the coefficients of the
remainder of \((-1)^d(\rho^\ell-\rho)\) on division by \(\overline A\).
The reduced equations therefore say that \(\overline A\) divides
\(\rho^\ell-\rho\). Now \(\rho^\ell-\rho=\prod_{s\in\FF_\ell}(\rho-s)\),
and polynomials over a field factor uniquely. Over any field that
contains \(\FF_\ell\), in particular over \(\Bbbk\), the monic divisors
of degree \(d\) are therefore the \(\binom{\ell}{d}\) polynomials
\eqref{eq:prime-labels}, and their coefficients lie in \(\FF_\ell\).
Each of them is a simple solution of the reduced equations. Write
\(\rho^\ell-\rho=\overline A\,\overline C\). For a variation \(h\) with
\(\deg h<d\), the derivative of the remainder is, up to the nonzero
scalar \((-1)^d\),
\begin{equation}\label{eq:prime-jacobian}
 h\longmapsto-\operatorname{rem}_{\overline A}(\overline Ch),
\end{equation}
because the dividend in \eqref{eq:prime-collapse} does not vary with
\(A\). Since \(\rho^\ell-\rho\) has no repeated factor, \(\overline A\)
and \(\overline C\) have no common factor. If \(\overline A\) divides
\(\overline Ch\), it therefore divides \(h\), and then \(h=0\) because
\(\deg h<d\). The linear map \eqref{eq:prime-jacobian} is thus
injective, hence invertible.

\emph{Step 3: integrality and rigidity.} A solution could have \(\ell\)
in a denominator, and then it could not be reduced. We show that this
does not happen. The argument is a dominant balance, with \(\ell\) in
the role of the small parameter: a solution that is large is rescaled to
size one, and the leading-order problem, which is the continuous one,
then leads to a contradiction. For \(\lambda\ne0\) and a monic
polynomial \(A\) of degree \(d\), put
\(A_\lambda(\rho)=\lambda^{-d}A(\lambda\rho)\). This is the monic
polynomial whose roots are those of \(A\) divided by \(\lambda\), and
its coefficients are \(\alpha'_j=\alpha_j/\lambda^j\). A polynomial
\(T\) in the \(\alpha_j\) and \(\rho\) that is weighted-homogeneous of
weight \(w\) satisfies
\(T(\alpha;\lambda\rho)=\lambda^{w}T(\alpha';\rho)\).

Let \(A\) be a solution of the matching equations with coefficients in
\(K\), and suppose that some root of \(A\) is not \(\ell\)-integral. Let
\(\lambda\) be a root of least valuation, so that
\(\theta=\operatorname{ord}(\lambda)<0\), and put \(B=A_\lambda\). The
roots of \(B\) are \(\ell\)-integral, hence so are its coefficients
\(\alpha'_j\), by \textup{(V)}. Decompose
\(\ell S_A=\sum_{w\le\ell}T^{(w)}\) into its parts of weight \(w\). They
have coefficients in \(\ZZ_{(\ell)}\) by Step~1, and
\(T^{(\ell)}=\ell H_A\). Since \(A\) divides \(S_A\) and
\(A(\lambda)=0\), we have \(S_A(\lambda)=0\). Dividing this equation by
\(\lambda^{\ell}\) gives
\[
 \ell H_B(1)=-\sum_{w<\ell}\lambda^{w-\ell}\,T^{(w)}(\alpha';1).
\]
Each term on the right has positive valuation, because
\(\operatorname{ord}(\lambda^{w-\ell})=(w-\ell)\theta>0\) and
\(T^{(w)}(\alpha';1)\) is \(\ell\)-integral. By
\eqref{eq:valuation-rules} the right side reduces to zero. The left side
reduces to \((-1)^d\ne0\), by \eqref{eq:prime-continuous} evaluated at
\(\alpha=\alpha'\) and \(\rho=1\). This contradiction shows that every
root of \(A\) is \(\ell\)-integral, and so is every coefficient. The
solution can be reduced, and by Step~2 its reduction is one of the
polynomials \eqref{eq:prime-labels}.

Rigidity is proved in the same way. Suppose that \(A\ne\rho^d\) is monic
of degree \(d\), with complex coefficients, and divides \(H_A\). Since
\(H_{A_\lambda}(\rho)=\lambda^{-\ell}H_A(\lambda\rho)\), every
\(A_\lambda\) has the same property, so we may assume that one of the
coefficients \(\alpha_j\) equals \(1\). By \textup{(N)}, applied to the
equations \(\remA H_A=0\) together with this normalization, we may
assume that the coefficients of \(A\) lie in \(\overline\QQ\subset K\).
Let \(\lambda\) be a root of \(A\) of least valuation; it is not zero,
because \(A\ne\rho^d\). As before, \(B=A_\lambda\) has \(\ell\)-integral
coefficients. Since \(A\) divides \(H_A\) and \(A(\lambda)=0\), we have
\(0=\ell H_A(\lambda)=\lambda^{\ell}\,\ell H_B(1)\). But \(\ell H_B(1)\)
reduces to \((-1)^d\ne0\), by \eqref{eq:prime-continuous}. This
contradiction proves~(c).

\emph{Step 4: lifting.} The system \((\ell R_r)_{0\le r<d}\) has
coefficients in \(\ZZ_{(\ell)}\), and by Step~2 each of the
\(\binom{\ell}{d}\) reduced solutions lies in \(\FF_\ell^d\) and is
simple. By Hensel's lemma \textup{(H)}, each of them is the reduction of
exactly one solution in \(\mathcal O^d\). By Step~3 every solution in
\(K^d\) lies in \(\mathcal O^d\), and its reduction is one of the
reduced solutions. Hence reduction modulo \(\ell\) is a bijection
between the solutions in \(K^d\) and the reduced solutions. There are
exactly \(\binom{\ell}{d}\) solutions in \(K^d\), and the Jacobian
determinant is nonzero at each of them, because its reduction is
nonzero. By \textup{(T)}, the matching equations have the same
\(\binom{\ell}{d}\) solutions over \(\CC\), with coordinates in
\(\overline\QQ\), all simple. This proves (a) and (b); part~(c) was
proved in Step~3.
\end{proof}

\begin{remark}[Formal verification]\label{rem:lean}
Parts (a) and (c) of Theorem~\ref{thm:prime} have been verified in
Lean~4 against the Mathlib library, both for the matching equations and
for the pairs \((P,v)\) themselves (Appendix~\ref{app:verification}).
The formal proof follows the route above, and like the proof above it
normalizes the roots of \(A\).
\end{remark}

\subsection{Reading the labels: pulses and fronts}\label{sec:labels}

As \(z\to+\infty\) we have \(Q\to1\) and \(D^jQ\to0\) for \(j\ge1\), so
\(F(D)Q\to F(0)\) for every polynomial \(F\). The limits of \(v_A\) as
\(z\to\mp\infty\) are therefore \(0\) and \(s_pA(0)\): pairs with
\(A(0)=0\) are pulses, and the others are fronts. The pulses are needed
separately, because the elliptic pairs of
Section~\ref{sec:elliptic-counting} degenerate onto them.

Letting \(z\to+\infty\) in \eqref{eq:discrete-square} gives
\(S_A(0)=-A(0)^2\). Since \(S_A(0)=A(0)C_A(0)+R_0\) by
\eqref{eq:remainder}, this gives the polynomial identity
\begin{equation}\label{eq:constant-remainder}
 R_0=-A(0)\bigl(A(0)+C_A(0)\bigr).
\end{equation}
On the hyperplane \(\alpha_d=A(0)=0\) the equation \(R_0=0\) thus holds
identically, and the pulses are the solutions of the \emph{pulse equations}
\begin{equation}\label{eq:pulse-equations}
 R_1=\cdots=R_{d-1}=0\qquad\text{in the unknowns }
 \alpha_1,\ldots,\alpha_{d-1},\quad\alpha_d=0 .
\end{equation}
We write \(\mathcal Z_p\) for the set of their solutions; multiplicities
of points of \(\mathcal Z_p\) refer to the pulse equations.

\begin{proposition}[Pulses and fronts]\label{cor:endpoints}
Let \(\ell=2p-1\) be prime. A pair is a pulse if and only if its label
contains \(0\). Consequently exactly \(S_p:=\binom{2p-2}{p-2}\)
of the \(N_p\) pairs are pulses, and the remaining
\(\binom{2p-2}{p-1}\) are fronts. Every pulse is a simple solution of
the pulse equations \eqref{eq:pulse-equations}, and
\(P_A(0)\ne0\) at every pulse (these two facts are used in
Section~\ref{sec:elliptic-counting}). These are counts over \(\CC\).
\end{proposition}

In Table~\ref{tab:ks} the pulses are the four pairs with \(b=0\), and
their labels are the four that contain \(0\).

\begin{proof}
Let \(A\) be a solution with label \(L\). If \(A(0)=0\), then
\(\overline A(0)=0\), so \(0\in L\). Conversely, let \(0\in L\). At the
solution \(A\) we have \(\ell S_A=A\cdot\ell C_A\). By \textup{(V)} and
\eqref{eq:prime-collapse}, \(\ell C_A\) is \(\ell\)-integral and reduces
to \((-1)^d\overline C\), where
\(\overline C=(\rho^\ell-\rho)/\overline A\). Since \(\rho\) divides
\(\overline A\) and \(\rho^\ell-\rho\) has no repeated factor,
\(\overline C(0)\ne0\). The number \(\ell A(0)\) reduces to zero. Hence
\(\ell\bigl(A(0)+C_A(0)\bigr)\) reduces to \((-1)^d\overline C(0)\ne0\).
So \(A(0)+C_A(0)\ne0\), and \(R_0=0\) gives \(A(0)=0\), by
\eqref{eq:constant-remainder}: the pair is a pulse. The number of
\(d\)-element subsets of \(\FF_\ell\) that contain \(0\) is
\(\binom{\ell-1}{d-1}=S_p\).

At a pulse, \eqref{eq:constant-remainder} shows that the row of the
Jacobian matrix of the matching equations that belongs to \(R_0\) has a
single entry that can be different from zero, namely \(-C_A(0)\) in the
column of \(\alpha_d\). The Jacobian matrix is invertible, by
Theorem~\ref{thm:prime}(a). Hence \(C_A(0)\ne0\), so that
\(P_A(0)=s_pC_A(0)/2\ne0\), and the minor obtained by deleting that row
and that column is invertible. This minor is the Jacobian matrix of the
pulse equations. (For \(d=1\) there are no pulse equations, and only
\(C_A(0)\ne0\) is asserted.)
\end{proof}

\section{The count of elliptic pairs}\label{sec:elliptic-counting}

We now count the elliptic pairs of Section~\ref{sec:setting} on a fixed
lattice \(\Lambda\subset\CC\), and put \(g_i=g_i(\Lambda)\). The pole is
at zero, and the constant term of the profile is retained.
Neither reflection nor a lattice automorphism is
factored out: pairs related in this way are counted separately (so that
\(M_3=8\) below, and not fewer). Two lattices differ by a rescaling
exactly when they have the same invariant
\(j(\Lambda)=1728g_2^3/(g_2^3-27g_3^2)\).

We write elliptic functions through the Weierstrass relations
\begin{equation}\label{eq:weierstrass}
 Y^2=4X^3-g_2X-g_3,\qquad
 \mathscr D X=Y,\qquad
 \mathscr D Y=6X^2-\frac{g_2}{2}.
\end{equation}
For \(g_2^3\ne27g_3^2\), that is, on a lattice, these relations are
realized by \(X=\wp(z;g_2,g_3)\), \(Y=\wp'(z;g_2,g_3)\) and
\(\mathscr D=\mathrm d/\mathrm dz\); see \citet[Chapter~23]{DLMF}.
We use \(X\) and \(Y\) because the formulas below must keep their
meaning for the degenerate parameters \(g_2^3=27g_3^2\), which belong to
no lattice. Using \eqref{eq:weierstrass}, every polynomial in \(X\) and
\(Y\) reduces to a unique normal form \(F(X)+YG(X)\). We say that
\(X^kY^\epsilon\), \(\epsilon\in\{0,1\}\), has pole order
\(2k+3\epsilon\); on a lattice this is its pole order at zero.

The representation of the profile is the classical one. An elliptic
function with one pole per period has zero residue there.
The functions \(1,\wp,\wp',\ldots,\wp^{(p-2)}\) have distinct
principal parts and span the space of elliptic functions whose only
pole is at zero, of order at most \(p\).
Subtracting the indicated principal part from a one-pole solution
leaves an entire elliptic function, hence a constant.
Since \(D^{p-2}\wp=(-1)^{p}(p-1)!\,z^{-p}+O(1)\), the leading
coefficient \eqref{eq:leadingpole} shows that every profile has the form
\begin{equation}\label{eq:elliptic-enumeration-profile}
 V=-s_p\mathscr D^{p-2}X
       +\sum_{j=0}^{p-3}u_j\mathscr D^jX+b.
\end{equation}
Its only pole on \(\CC/\Lambda\) is zero, of order \(p\).
In particular \(\Lambda\) is its full period lattice.
The matching equations \(\mathcal E_p(g_2,g_3)\) are obtained by formal
substitution of \eqref{eq:elliptic-enumeration-profile} in
\eqref{eq:main}, using only \eqref{eq:weierstrass}: the left side is
reduced to its normal form, and the coefficients of \(F\) and \(G\) are
set equal to zero. The unknowns are \(a_0,\ldots,a_{p-1}\),
\(u_0,\ldots,u_{p-3}\) and \(b\). The equations are polynomials in the
unknowns and in \(g_2,g_3\), and they are defined for every
\((g_2,g_3)\in\CC^2\), the degenerate parameters included; on a lattice
their solutions are the elliptic pairs. The symbol
\(\mathcal E_p(g_2,g_3)\) also denotes the set of solutions.
Multiplicity and simplicity are understood as
in Section~\ref{sec:counting}.

\begin{theorem}[Fixed-lattice elliptic enumeration]\label{thm:elliptic-count}
Let \(2p-1\) be prime.
\begin{enumerate}
\item[\textup{(a)}] On every lattice the elliptic pairs are
finite in number, and counted with multiplicity there are
\begin{equation}\label{eq:elliptic-length}
 M_p:=2\binom{2p-2}{p-2}=2S_p
\end{equation}
of them.
\item[\textup{(b)}] There is a proper algebraic subset \(\mathcal N\) of
the \((g_2,g_3)\)-plane outside which all of them are simple. The
lattices in \(\mathcal N\) are those whose invariant \(j(\Lambda)\) lies
in a finite set, which depends on \(p\). Thus a generic lattice, that
is, a lattice whose invariant avoids these values, admits exactly
\(M_p\) distinct elliptic pairs.
\end{enumerate}
\end{theorem}

\paragraph{The smallest case.}
For \(p=2\) the profile is \(V=-12X+b\), and the left side of
\eqref{eq:main} reduces to
\[
 -12a_1Y-12(a_0+b)X+6g_2+a_0b+\frac12b^2.
\]
Consequently the complete elliptic system is
\begin{equation}\label{eq:kdvb-elliptic}
 a_1=0,\qquad
 v=-12\wp-a_0,\qquad a_0^2=12g_2.
\end{equation}
There is no one-pole elliptic member with nonzero Burgers coefficient:
both elliptic pairs belong to the KdV equation, in agreement with
\(M_2=2\). They are the two backgrounds \eqref{eq:shift} of one
cnoidal wave. For \(g_2\ne0\) they are distinct and simple.
For \(g_2=0\), \(g_3\ne0\), they coincide, and the lattice carries one
pair of multiplicity two. Here \(\mathcal N\) is the line \(g_2=0\),
and the finite set of (b) is \(\{0\}\). At order three it is
\(\{0,-300\}\) (Section~\ref{sec:ks}).

\paragraph{No escape to infinity.}
The number \(M_p\) is twice the number \(S_p\) of pulses of
Proposition~\ref{cor:endpoints}, and the proof shows why.
When the parameters of a square polynomial system (as many equations as
unknowns) vary, its solutions may merge or separate. As long as they
remain finite in number, their number, counted with multiplicity, can
change only if a solution moves off to infinity. A solution that moves
off to infinity, rescaled to size one, converges to a nonzero solution
of the leading parts of the equations; this is the rescaling argument
of Step~3 in the proof of Theorem~\ref{thm:prime}, with the ordinary
absolute value in place of the \(\ell\)-adic one. Here the leading
parts are the equations with both periods infinite
(\(\wp\to z^{-2}\); the cusp \(g_2=g_3=0\)), and by rigidity that
rational problem has only the trivial solution. So nothing escapes, and
the number of pairs, counted with multiplicity, is the same for all
\((g_2,g_3)\). We find it with one period infinite (\(\wp\) becomes a
hyperbolic function; the node): there every elliptic pair becomes a
pulse on one of its two constant backgrounds, and every pulse on either
background arises in this way exactly once. The pulses are simple, and
simplicity spreads from the node to a generic lattice.

\begin{lemma}[Conservation of number]\label{lem:conservation}
Let \(E_1,\ldots,E_n\) be polynomials in unknowns
\(x=(x_1,\ldots,x_n)\) and parameters \(t=(t_1,\ldots,t_m)\),
weighted-homogeneous for some positive integer weights of the \(x_j\)
and the \(t_k\). Suppose that the leading system \(E_i(x;0)=0\),
\(1\le i\le n\), has no solution in \(\CC^n\) other than \(x=0\). Then:
\begin{enumerate}
\item[\textup{(i)}] for every \(t\) the system \(E(x;t)=0\) has
finitely many solutions, and they stay bounded while \(t\) stays
bounded;
\item[\textup{(ii)}] the number of solutions, counted with
multiplicity, is the same for every \(t\in\CC^m\);
\item[\textup{(iii)}] the \(t\) for which some solution is not simple
form an algebraic subset of \(\CC^m\), that is, the common zero set of
finitely many polynomials in \(t\).
\end{enumerate}
\end{lemma}

The proof of the lemma uses more algebra than the rest of the paper and
can be skipped; Steps~1--5 below use only its statement.

\begin{proof}
Since the leading system has only the zero solution, the polynomials
in \(x\) modulo its equations form a finite-dimensional space
\citep[Chapter~5, \S3]{CoxLittleOSheaIVA}, spanned by monomials
\(m_1=1,\ldots,m_r\). Every polynomial \(f\) in \((x,t)\) is then,
modulo the equations \(E_i\), a combination \(\sum_ic_i(t)m_i\) with
polynomials \(c_i\) in \(t\). This is proved by induction on the
weight of \(f\), which may be taken weighted-homogeneous: write
\(f=f(x,0)+\sum_kt_kf_k\), write \(f(x,0)\) as a combination of the
\(m_i\) plus \(\sum_ih_iE_i(x;0)\) with weighted-homogeneous \(h_i\),
and replace each \(E_i(x;0)\) by \(E_i\) minus its terms that contain
some \(t_k\). What remains is a sum of products of a \(t_k\) with a
polynomial of lower weight, because \(t_k\) has positive weight.
In particular \(x_jm_i=\sum_lM^{(j)}_{il}(t)m_l\) modulo the equations.
Multiplying this vector identity by the adjugate of the matrix
\(x_j-M^{(j)}(t)\) shows that \(\det\bigl(x_j-M^{(j)}(t)\bigr)\,m_i\) is
a combination of the \(E_i\) for every \(i\) (the Cayley--Hamilton
argument); for \(m_1=1\) this is the determinant itself.
So at every solution \(x_j\) is a root of a
monic polynomial of degree \(r\) whose coefficients are polynomials in
\(t\). This proves~(i).

For (ii), join two parameter values by a segment. By (i) some ball
\(|x|<\varrho\) contains all solutions for all \(t\) on the segment, so
no solution crosses the sphere \(|x|=\varrho\). The number of zeros of
a holomorphic map in a ball, counted with multiplicity, when they are
finite in number and none lies on the bounding sphere, is the degree
of \(E/|E|\) on that sphere, and it does not change under a
deformation without zeros on the sphere
\citep[Appendix~B]{Milnor1968},
\citep[Sections~5.2 and~5.4]{ArnoldGuseinZadeVarchenko1985}; this is
Rouch\'e's theorem in several variables.

For (iii), the set in question is the projection to \(\CC^m\) of the
algebraic set on which \(E_1,\ldots,E_n\) and the Jacobian determinant
\(\det(\partial E_i/\partial x_j)\) vanish.
The monic polynomials of (i) are combinations of the \(E_i\). Eliminate
the unknowns one at a time: at each step the polynomial of (i) for the
unknown being eliminated has leading coefficient \(1\), so by the
Extension Theorem of elimination theory
\citep[Chapter~3, \S1]{CoxLittleOSheaIVA} every solution of the
eliminated equations extends. The projection is therefore the common
zero set of the polynomials in \(t\) alone that are combinations of the
\(E_i\) and the Jacobian determinant.
\end{proof}

Three operations do not change multiplicities, and the proof of the
theorem uses them freely: an invertible affine change of the unknowns;
invertible linear combinations of the equations with constant
coefficients; and the elimination of an unknown that one of the
equations gives as a polynomial in the others. Each of them leaves
unchanged, up to an isomorphism, the space of power series modulo the
combinations of the equations, whose dimension is the multiplicity
(Section~\ref{sec:counting}).

\begin{proof}[Proof of Theorem~\ref{thm:elliptic-count}]
\emph{Step 1: eliminating the operator.}
The coefficient at pole order \(2p\) vanishes by the choice of
\(s_p\).
The coefficients at orders \(2p-1,\ldots,p\) determine
\(a_{p-1},\ldots,a_0\) successively.
At each step the coefficient of the new \(a_j\) is the nonzero
rational leading coefficient of \(D^jV\).
This is a polynomial elimination: each \(a_j\) is obtained as a
polynomial in the remaining unknowns and in \(g_2,g_3\), and
multiplicities do not change.
The residual has pole order at most \(p-1\) and zero residue, since
no normal form has pole order one. It lies in the span of \(1\) and the
\(\mathscr D^jX\), \(0\le j\le p-3\).
Hence \(p-1\) equations \(E_1=\cdots=E_{p-1}=0\) remain in the \(p-1\)
profile coefficients \(u_j,b\): the coefficients of the residual in this
basis.

\emph{Step 2: weights.}
Give \(D,X,g_2,g_3\) weights \(1,2,4,6\), respectively, and give
the profile weight \(p\); this is the usual homogeneity of \(\wp\)
under a rescaling of \(z\).
Then \(u_j\) has weight \(p-j-2\), and \(b\) has weight \(p\).
After the triangular elimination, the residual principal-part
coefficients and the residual constant, that is, the \(E_i\), are
weighted-homogeneous polynomials in the \(u_j\), \(b\), \(g_2\),
\(g_3\), of weights
\begin{equation}\label{eq:elliptic-weights}
 \begin{array}{c|c}
 \text{variables}&1,2,\ldots,p-2,p\\
 \text{equations}&p+1,p+2,\ldots,2p-2,2p.
 \end{array}
\end{equation}
The B\'ezout number of this table is
\((p+1)\cdots(2p-2)(2p)/\bigl((p-2)!\,p\bigr)=2\binom{2p-2}{p-2}=M_p\)
(at \(p=3\): \(4\cdot6/(1\cdot3)=8\)), in agreement with the count
below (companion note).
The normal forms and the eliminations are polynomial in
\(g_2,g_3\), so these statements extend to the degenerate parameters.
For fixed \(g_2,g_3\), the highest weighted parts of the equations in
the profile variables are obtained by setting \(g_2=g_3=0\), since
\(g_2\) and \(g_3\) carry positive weight.

\emph{Step 3: the cusp.}
At the cusp \(g_2=g_3=0\) the relations \eqref{eq:weierstrass} are
satisfied by \(X=z^{-2}\), \(Y=-2z^{-3}\) and
\(\mathscr D=\mathrm d/\mathrm dz\), so a solution of the cusp system
gives a nonconstant rational solution \(V\) of \eqref{eq:main} whose
only pole is at zero. By Lemma~\ref{lem:laplace}, \(V\) vanishes at
infinity, so that
\(b=0\); moreover \(a_0=P(0)=0\), and \(V=(-1)^{p-1}s_pA^\sharp\) for a
monic \(A\) of degree \(p-1\) that divides \(H_A\).
By rigidity, which holds by Theorem~\ref{thm:prime}(c), \(A=\rho^{p-1}\).
Thus \(V=c_\ast z^{-p}\), and all the variables in
\eqref{eq:elliptic-enumeration-profile} vanish, because the functions
\(D^jz^{-2}\) have different pole orders.
The leading equations have no nonzero common zero.

\emph{Step 4: conservation of number.}
The residual equations \(E_1,\ldots,E_{p-1}\) of Step~1
are polynomials in the \(p-1\) profile unknowns and in \(g_2,g_3\),
weighted-homogeneous for the weights of Step~2, and by Step~3 their
leading system, which is the cusp system, has only the zero solution.
Lemma~\ref{lem:conservation} applies, with \(t=(g_2,g_3)\). By (i),
every system \(\mathcal E_p(g_2,g_3)\) has finitely many solutions,
which proves the finiteness in (a). By (ii), their number, counted with
multiplicity, is the same at every point of
the \((g_2,g_3)\)-plane, the degenerate parameters included; at the
cusp all of this multiplicity sits
at the single solution where every unknown vanishes. By (iii), the
parameters at which some solution is not simple form an algebraic
subset \(\mathcal N\) of the plane.

\emph{Step 5: the node.}
To find the common number, and to show that \(\mathcal N\) is not the
whole plane, consider the nodal parameters
\begin{equation}\label{eq:nodal-parameters}
 g_2=\frac1{12},\qquad g_3=-\frac1{216}.
\end{equation}
For these parameters the relations \eqref{eq:weierstrass} are
satisfied by
\[
 X=Q^2-Q+\frac1{12},\qquad
 Y=(2Q-1)Q(1-Q),\qquad \mathscr D=D\ \text{ with }\ DQ=Q(1-Q);
\]
this is the degenerate \(\wp\)-function with the single period
\(2\pi i\), with its pole moved to \(z=i\pi\), written in the logistic
coordinate. A monomial \(X^kY^\epsilon\) becomes a polynomial in
\(Q\) of degree \(2k+3\epsilon\), which is also its pole order. These
degrees are all different, so the monomials remain linearly
independent, and the matching equations for the nodal parameters are
exactly the conditions for \eqref{eq:main} to hold as an identity
between polynomials in \(Q\).
Since \(X=\frac1{12}-DQ\) and \(\mathscr D^jX=-D^{j+1}Q\) for
\(j\ge1\), the profile \eqref{eq:elliptic-enumeration-profile}
becomes, with \(u_{p-2}=-s_p\),
\[
 V=\beta+v_A,\qquad
 A(\rho)=-\frac1{s_p}\sum_{j=0}^{p-2}u_j\rho^{j+1},\qquad
 \beta=b+\frac{u_0}{12}.
\]
Here \(A\) is monic of degree \(d=p-1\) with \(A(0)=0\), and \(\beta\)
is the common limit of \(V\) at the two ends. With the coefficients
\(\alpha_j\) of \(A\) as in \eqref{eq:normalized-A}, the passage from
\((u_0,\ldots,u_{p-3},b)\) to
\((\alpha_1,\ldots,\alpha_{d-1},\beta)\) is an invertible affine
change of unknowns with rational coefficients.

In these unknowns the constant term of \(P(D)V+V^2/2\) is
\(P(0)\beta+\beta^2/2\), and the remaining terms are
\((P+\beta)(D)v_A+v_A^2/2\). By the proof of
Proposition~\ref{prop:divisibility} the latter equals
\((s_p/2)\bigl[2(P+\beta)A-s_pS_A\bigr](D)Q\), and the elimination of
Step~1 is the division of \(S_A\) by \(A\): it
gives \(P+\beta=s_pC_A/2\), a polynomial in the unknowns, which we
still denote by \(P_A\), and it leaves \(-(s_p^2/2)R_A(D)Q\).
Since \(A(0)=0\), the coefficient \(R_0\) vanishes identically, by
\eqref{eq:constant-remainder}. With \(P(0)=P_A(0)-\beta\), the nodal
matching equations are therefore invertible linear combinations, with
constant coefficients, of
\begin{equation}\label{eq:nodal-matching}
 R_1=\cdots=R_{d-1}=0,\qquad
 \beta\bigl(\beta-2P_A(0)\bigr)=0,
\end{equation}
together with \(P=P_A-\beta\). These are the pulse equations
\eqref{eq:pulse-equations} and one quadratic equation for the
background, and the two systems have the same solutions with the same
multiplicities.

By Proposition~\ref{cor:endpoints} the pulse equations have exactly
\(S_p\) solutions, all simple, with \(P_A(0)\ne0\) at each of them.
For each of them the quadratic has the two simple roots \(\beta=0\)
and \(\beta=2P_A(0)\), and the Jacobian matrix of
\eqref{eq:nodal-matching} is block triangular. For the nodal
parameters the system thus has exactly \(2S_p\) solutions, all simple.
By Step~4 the number of solutions counted with multiplicity is
\(2S_p=M_p\) for all \((g_2,g_3)\). This proves (a). Moreover the
nodal parameters do not belong to \(\mathcal N\). So \(\mathcal N\) is
a proper algebraic subset of the plane, and outside \(\mathcal N\)
every solution is simple. A nonzero polynomial cannot vanish on a
nonempty open subset of the plane, so \(\mathcal N\) does not contain
all the lattices. It contains only those with finitely many values of
\(j\), as we now show.

The equations \(E_i\) and their Jacobian determinant are
weighted-homogeneous, so \(\mathcal N\) is invariant under
\((g_2,g_3)\mapsto(\lambda^4g_2,\lambda^6g_3)\), \(\lambda\ne0\). On
lattices this is the rescaling \(\Lambda\mapsto\lambda^{-1}\Lambda\),
which maps pairs to pairs and simple pairs to simple pairs. If a
polynomial vanishes on \(\mathcal N\), then so does each of its
weighted-homogeneous components. Since \(\mathcal N\) is a proper
algebraic subset, it therefore lies in the zero set of a nonzero
polynomial that is a monomial in \(g_2,g_3\) times a product
\(\prod_i(g_2^3-c_ig_3^2)\). On the lattices this zero
set consists of finitely many values of \(j\): \(j=0\) on \(g_2=0\),
\(j=1728\) on \(g_3=0\), and \(j=1728c_i/(c_i-27)\) on
\(g_2^3=c_ig_3^2\). Two lattices with the same \(j\) differ by a
rescaling, so \(\mathcal N\) contains both or neither. Hence the
lattices in \(\mathcal N\) are exactly those whose invariant
\(j(\Lambda)\) lies in a finite set, which depends on \(p\).
This proves (b).
\end{proof}

\paragraph{Two backgrounds.}
The factor two in \(M_p\) retains the two choices of background in
\eqref{eq:shift}. That transformation preserves the evolution equation
and the lattice; it changes the operator \(P\) and the speed. For the
nodal parameters it exchanges the two roots of the quadratic in
\eqref{eq:nodal-matching}. Outside \(\mathcal N\) the two backgrounds of
every wave are distinct, because a solution with \(a_0=0\) is never
simple: Step~1 gives \(a_0+b\) as a polynomial in the \(u_j\) and
\(g_2,g_3\), and \(b\) enters the residual equations only through the
constant one, which gives \(a_0^2\) as another such polynomial; so the
Jacobian determinant contains the factor \(a_0\). On a generic lattice
the \(M_p\) pairs therefore represent \(M_p/2\) elliptic waves, up to
translations and Galilean shifts, each counted with both of its
backgrounds. On special lattices the two backgrounds may
coalesce, as \eqref{eq:kdvb-elliptic} shows for \(g_2=0\) and
Section~\ref{sec:ks} shows at order three.

The elliptic pairs of purely dispersive equations, \(2F_p\) of the
\(M_p\), are counted in Section~\ref{sec:reflection}
(Corollary~\ref{cor:elliptic-symmetric}), after the symmetry of their
profiles has been established; at order two they are both pairs of
\eqref{eq:kdvb-elliptic}, and \(M_2=2F_2=2\).

\section{Which orders are settled}\label{sec:further}

Each formula needs two facts at a given \(p\). The first is the
rigidity condition \((\mathrm R_{p-1})\) of Section~\ref{sec:counting}.
It gives \(N_p\), \(M_p\), \(F_p\) and \(2F_p\) as counts with
multiplicity: the multiplicities of the pairs add up to these numbers.
The second is that every rational-exponential pair is simple, that is,
\(\Xp\) is reduced. Then the pairs are distinct, and the elliptic pairs
are distinct on a generic lattice. On a special lattice several
elliptic pairs can merge, while the count with multiplicity stays the
same; Section~\ref{sec:ks} shows this at order three.
Table~\ref{tab:status} records what is known about each fact.

\begin{table}[htbp]
\centering
\caption{What is known. For elliptic pairs, distinctness is asserted
only on a generic lattice.}
\label{tab:status}
\begin{tabular}{@{}p{0.36\linewidth}p{0.58\linewidth}@{}}
\toprule
Statement&Known for\\
\midrule
Counts with multiplicity&\(2p-1\) a prime or a prime power (proved);
every \(p\le73\) (exact certificate); every \(p\le573\) (finite-field
computations)\\[2pt]
All \(N_p\) pairs distinct&\(2p-1\) prime (proved); \(p=5,8,11\) (exact
computations); hence every \(p\le12\)\\[2pt]
All \(F_p\) pairs of the pure subfamily distinct&whenever all \(N_p\)
pairs are distinct; every \(p\le14\) (exact computations)\\[2pt]
Both statements, every \(p\)&conjectured\\
\bottomrule
\end{tabular}
\end{table}

The prime-index argument excluded solutions at infinity by reducing
the continuous convolution modulo \(\ell\). The same argument proves
\eqref{eq:length} whenever \(2p-1=\ell^e\) is an odd prime power
(Appendix~\ref{app:verification}). It gives the count with
multiplicity, but it does not show that the pairs are distinct.

More generally, the exact arithmetic certificates supplied with the
paper establish \((\mathrm R_d)\) for every \(d\le72\), hence
\begin{equation}\label{eq:coverage}
 \len\Xp=\binom{2p-1}{p-1},\qquad
 \len\mathcal E_p(g_2,g_3)=2\binom{2p-2}{p-2}
 \qquad(2\le p\le73).
\end{equation}
In the same range, Theorem~\ref{thm:symmetric} and
Corollary~\ref{cor:elliptic-symmetric} give
\(\len\Xp^{\mathrm{sym}}=F_p\) and, for even \(p\),
\(\len\mathcal E_p^{\mathrm{sym}}(g_2,g_3)=2F_p\).
Together with the ancillary proof, the certificates also cover infinite
families of higher orders.
Finite-field Gr\"obner computations, checked by independent pipelines,
establish \((\mathrm R_d)\) for every \(d\le572\); they extend
\eqref{eq:coverage} and the two symmetric lengths, as counts with
multiplicity, to every \(2\le p\le573\) (repository, directory
\texttt{certificates/rigidity\_572/}).
The certificates hold for all complex coefficients, not only real or
rational ones. The ancillary proof and the full certificate are described
in Appendix~\ref{app:verification}.

Equation \eqref{eq:coverage} is a statement about counts with
multiplicity only. At composite indices, exact computations show that
all pairs are simple at \(p=5\), \(p=8\) and \(p=11\)
(Appendix~\ref{app:verification}). With Theorem~\ref{thm:prime}, all
pairs are therefore distinct for every \(p\le12\), and
Theorem~\ref{thm:elliptic-count} gives \(M_p\) distinct elliptic pairs
on a generic lattice at these orders. Inside the two pure subfamilies
the computation is much smaller, and it shows that the \(F_p\) pairs
are distinct for every \(p\le14\), which includes the composite indices
\(9\), \(15\), \(21\), \(25\) and \(27\).

\section{What the count means for a single equation}\label{sec:fixed}

The count lets the operator vary. We now read it for one evolution
equation \eqref{eq:pde} with fixed complex coefficients. The equation is
invariant under translations and under Galilean shifts.

\begin{proposition}[The waves of a fixed equation]\label{prop:fixed}
Let \(\kappa\ne0\). Up to a translation and a Galilean shift, the
travelling waves of \eqref{eq:pde} that are rational in
\(e^{\kappa(x-ct)}\), with one pole modulo
\((2\pi i/\kappa)\ZZ\), are
\begin{equation}\label{eq:fixed-wave}
 u=b_p\kappa^p\,v_A\bigl(\kappa(x-ct)\bigr),\qquad
 c=-b_p\kappa^pP_A(0),
\end{equation}
where \(A\) ranges over the points of \(\Xp\) that satisfy
\begin{equation}\label{eq:fixed-conditions}
 \frac{b_j}{b_p}=\kappa^{p-j}\,[\rho^j]P_A,\qquad 1\le j\le p-1 .
\end{equation}
The data \((A,\kappa)\) and \((\iota A,-\kappa)\) describe the same
wave, read in the two directions (Section~\ref{sec:reflection}).
\end{proposition}

\begin{proof}
Such a wave has a finite limit as \(e^{\kappa(x-ct)}\to0\)
(Section~\ref{sec:systems}), and a Galilean shift makes this limit
zero. Then \(w_\ast=0\) in \eqref{eq:scaling}, and
\eqref{eq:coefficients} becomes \eqref{eq:fixed-conditions} together
with the formula for \(c\). After a translation \(v\) has the
normalization of Section~\ref{sec:setting}, so \(v=v_A\) and \(P=P_A\)
by Proposition~\ref{prop:divisibility}. The last assertion follows from
\eqref{eq:reflection-pair}: the two waves differ by the constant
\(b_p\kappa^ps_pA(0)\), which is a Galilean shift.
\end{proof}

Given a pair, \eqref{eq:fixed-conditions} are \(p-1\) conditions on the
\(p-1\) coefficient ratios \(b_j/b_p\) and the one unknown \(\kappa\):
a pair imposes \(p-2\) conditions on the equation and fixes the rate.
A pair with real coefficients gives a real wave of a real equation,
smooth on the real axis, exactly when \(\kappa\) is real.
Let \(\mathbf a(A)=([\rho^{p-1}]P_A,\ldots,[\rho^{1}]P_A)\) be the
vector of lower coefficients of \(P_A\), and let
\(\mathbf b=(b_{p-1}/b_p,\ldots,b_1/b_p)\) be the coefficient vector of
the equation. Condition \eqref{eq:fixed-conditions} says that
\(\mathbf b\) is obtained from \(\mathbf a(A)\) by multiplying its
\(k\)th entry by \(\kappa^k\), for some \(\kappa\ne0\). As \(\kappa\)
varies these vectors trace a curve, denoted
\(\CC^\ast\mathbf a(A)\), and rescaling \(x\) moves an equation along
it. This is the reason for fixing the rate: without it every pair would
be a one-parameter family. Each of the \(N_p\) pairs marks one curve,
and several pairs may mark the same curve.

\begin{corollary}[Solvable equations]\label{cor:solvable}
Assume that the matching equations have finitely many solutions.
This holds when \(2p-1\) is prime, and whenever rigidity holds
(Proposition~\ref{prop:bezout}), hence for every \(p\le573\)
(Section~\ref{sec:further}).
\begin{enumerate}
\item[\textup{(i)}] Equation \eqref{eq:pde} has a travelling
wave that is rational in an exponential, with one pole per period, if
and only if \(\mathbf b\) lies on one of the curves
\(\CC^\ast\mathbf a(A)\), \(A\in\Xp\). There are at most
\((N_p+F_p)/2\) such curves.
\item[\textup{(ii)}] If \(\mathbf b\ne0\), there are finitely many such
waves up to translations and Galilean shifts.
\item[\textup{(iii)}] If \(p\ge3\), the curves form a set of dimension
at most one in \(\CC^{p-1}\). A generic equation
\eqref{eq:pde} has no such wave.
\end{enumerate}
\end{corollary}

\begin{proof}
Part~(i) restates Proposition~\ref{prop:fixed}. For the bound,
\(\mathbf a(\iota A)=(-1)\cdot\mathbf a(A)\) by
\eqref{eq:reflection-pair}, so \(A\) and its mirror image \(\iota A\)
give the same curve. The number of classes \(\{A,\iota A\}\) is half
the number of points plus half the number of fixed points. The
solutions are finite in number, hence isolated, so these two numbers
are at most \(N_p\) and \(F_p\) (Proposition~\ref{prop:bezout},
Remark~\ref{rem:Fp-bound}). For~(ii), if \(\mathbf b\ne0\) and
\(\mathbf b\in\CC^\ast\mathbf a(A)\), some coordinate \([\rho^j]P_A\)
is nonzero, and \eqref{eq:fixed-conditions} leaves at most \(p-j\)
values of \(\kappa\). Part~(iii) is a dimension count.
\end{proof}

When \(\mathbf b=0\) the equation is invariant under the scaling itself,
and its waves come in one-parameter families of rates. At order two
there are two curves, and both carry waves: \(\mathbf b=0\) is the KdV
equation, with its solitons of all rates, and \(\mathbf b\ne0\) consists
of the equations with a nonzero Burgers coefficient, each with one
family of fronts, at the rate fixed by that coefficient
(Section~\ref{sec:kdvb}).

At \(p=3\), write \(P=D^3+BD^2+\mu D+a_0\) and use \((B,\mu)\) as the
coordinates on \(\CC^2\). Rescaling multiplies \(B\) by \(\kappa\) and
\(\mu\) by \(\kappa^2\), so the curves are \(B^2=\sigma^2\mu\) with the
origin removed.
For the ten pairs \(A=\rho^2+a\rho+b\) of Table~\ref{tab:ks},
\eqref{eq:ksoperator} gives \(B=4a\) and \(\mu=a^2+19b\), hence the
four classical values
\begin{equation}\label{eq:ks-sigma}
 \sigma^2=\frac{16a^2}{a^2+19b}\in
 \{0,\ 16,\ 144/47,\ 256/73\},
 \qquad\text{that is,}\qquad
 |\sigma|\in\{0,\ 4,\ 12/\sqrt{47},\ 16/\sqrt{73}\},
\end{equation}
which are those of \citet[Eq.~(39)]{Kudryashov1989}. In the
normalization of the Introduction, with \(b_1=b_3=1\) and
\(b_2=\sigma\), condition \eqref{eq:fixed-conditions} reads
\(\kappa^2=1/(a^2+19b)\) and \(\sigma=4a\kappa\).
Here the bound \((N_3+F_3)/2=6\) is not attained, because the two
pairs with \(a=0\) share \(\sigma=0\), and \((\pm1,0)\) and
\((\pm i,0)\) share \(\sigma^2=16\).
For coefficients off these curves, approximate solutions close to the
exact ones are constructed by \citet{KudryashovKochanov2012}.

For odd \(p\) the classification turns
Corollary~\ref{cor:solvable} into a statement about all meromorphic
travelling waves.

\begin{corollary}[Generic equations with \(p\) odd]
\label{cor:generic}
Let \(p\ge5\) be odd and assume rigidity \((\mathrm R_{p-1})\), which
holds when \(2p-1\) is prime and for every \(p\le573\). The coefficient
vectors \(\mathbf b\in\CC^{p-1}\) for which \eqref{eq:pde} has a
nonconstant meromorphic travelling wave lie in an algebraic subset of
dimension at most two. In particular a generic equation has none.
\end{corollary}

\begin{proof}
By Theorem~\ref{thm:complete}, such a wave is rational, rational in an
exponential, or elliptic, with one pole per period. Take
\(\kappa=1\) in \eqref{eq:scaling}, so that
\([\rho^j]P=b_j/b_p\) for \(1\le j\le p-1\).
The rational-exponential waves give a set of dimension at most one, by
Corollary~\ref{cor:solvable}. A rational wave is, after a translation,
of the form \eqref{eq:rational-system}. By Lemma~\ref{lem:laplace} it
comes from a monic \(A\) of degree \(p-1\) that divides \(H_A\). By
rigidity, \(A=\rho^{p-1}\), hence \(P=D^p\) and \(\mathbf b=0\).
For the elliptic waves, Steps~1--3 of the proof of
Theorem~\ref{thm:elliptic-count} use only rigidity, and with them
Lemma~\ref{lem:conservation}(i) shows that every system
\(\mathcal E_p(g_2,g_3)\) has finitely many solutions. The solutions
for all \((g_2,g_3)\) together therefore form an algebraic set of
dimension at most two, and so does the closure of its image under
\(P\mapsto\mathbf b\).
\end{proof}

At \(p=3\) the argument gives the exact answer: the set of these
\(\mathbf b=(B,\mu)\) is the union of the four curves
\(B^2=\sigma^2\mu\) of \eqref{eq:ks-sigma}. Indeed the elliptic waves
require \(B^2=16\mu\) (Section~\ref{sec:ks}), which is one of the four
curves, and the point \(\mathbf b=0\) of the rational wave lies on all
of them.

The same dimension count applies to the one-pole waves of every order
\(p\ge3\) for which rigidity holds: a generic equation has no
travelling wave that is rational, rational in an exponential, or
elliptic with one pole per period. For even \(p\), however,
Theorem~\ref{thm:complete} excludes other meromorphic waves only when
the equation is purely dispersive.

Can one equation have two different waves of this kind at the same
rate and speed? It cannot, and this holds for every \(P\), including
the mixed equations of odd order, to which Theorem~\ref{thm:complete}
does not apply: the argument only uses the leading coefficient of a
convolution.

\begin{proposition}[One profile per operator]\label{prop:injective}
Let \(A_1,A_2\) be monic polynomials of degree \(p-1\) satisfying
\(A_1\mid S_{A_1}\) and \(A_2\mid S_{A_2}\). If \(P_{A_1}=P_{A_2}\),
then \(A_1=A_2\).
Consequently a monic operator of degree \(p\) admits at most one
normalized rational-exponential one-pole profile.
Any count of distinct pairs with this normalization is therefore
also a count of distinct operators.
\end{proposition}

\begin{proof}
Suppose \(P_{A_1}=P_{A_2}=P\), and put \(h=A_1-A_2\).
If \(h\ne0\), its degree \(n\) is less than \(p-1\); write
\(f\ne0\) for its leading coefficient.
By \eqref{eq:operator-A} and the bilinearity and symmetry of the
convolution, \(S_{A_1}-S_{A_2}=h\star(A_1+A_2)=(2/s_p)Ph\).
The power-sum calculation in \eqref{eq:Sleading} gives the leading
coefficient \(2f\int_0^1t^n(1-t)^{p-1}\dd t\) on the left.
Since \(P\) is monic, the right side has leading coefficient
\((2/s_p)f\), and \(2/s_p\) is the beta integral of
\eqref{eq:Sleading}. Cancelling \(f\) would therefore give
\[
 2\int_0^1t^n(1-t)^{p-1}\dd t
   =\int_0^1t^{p-1}(1-t)^{p-1}\dd t.
\]
This is impossible: \(n<p-1\) implies \(t^n>t^{p-1}\) for \(0<t<1\), so
the left side exceeds twice the right side. Thus \(h=0\).
\end{proof}

Every \(P_A\) has a root among \(1,\ldots,p\). Let \(m\) be the least
positive integer with \(A(m)\ne0\); then \(m\le p\), because \(A\) has
degree \(p-1\). As \(e^z\to0\) the profile begins with
\(s_p(-1)^{m-1}A(m)e^{mz}\) and its square with \(e^{2mz}\), so the
coefficient of \(e^{mz}\) in \eqref{eq:main} gives \(P_A(m)=0\).
In physical terms, the wave leaves its background like \(e^{mz}\), so
\(m\) is a spatial eigenvalue of the linearization about the
background, that is, a root of \(P\); and \(m\) is one of
\(1,\ldots,p\).

\paragraph{Other equations with the same profile equation.}
The reduction of Section~\ref{sec:setting} uses only the travelling-wave
variable. The counts therefore apply to any equation
\(\mathcal L u+\mathcal M(u^2)=0\) with constant-coefficient operators,
mixed derivatives and several space variables included, whose
plane-wave reduction integrates to \eqref{eq:main}; the coefficients of
\(P\) then depend on the speed and the direction of the wave.
An example in three space variables is the equation of
\citet{KudryashovSinelshchikovThreeDim2012} for waves in a liquid with
gas bubbles, whose travelling-wave reduction is that of KdV--Burgers.
The repository lists such equations by profile operator, with their
reductions (directory \texttt{tools/wave\_dictionary/}).

\paragraph{Ansatz methods.}
Many exact travelling waves in the literature are produced by
substituting a finite expansion: the tanh and sech methods, the
exp-function and \((G'/G)\)-expansion methods, the simplest-equation
method with a Riccati equation \citep{Kudryashov2005Simplest}, and the
Jacobi and Weierstrass elliptic-function expansions
\citep{Kudryashov2009Errors}. In their standard form, with one phase
\(kx-\omega t\) and constant coefficients, the first group produces
functions that are rational in one exponential and the second group
elliptic functions of one argument; rational functions of \(z\) arise
as limits. With respect to its primitive period, an output with one
pole per period is, up to the symmetries of
Proposition~\ref{prop:fixed}, one of the waves \eqref{eq:fixed-wave},
one of the elliptic waves of Section~\ref{sec:elliptic-counting}, or
the rational wave; for the equations of Theorem~\ref{thm:complete}
every nonconstant output has one pole per period. The counts therefore
serve as a checklist: at a covered order, once the rate or the period
lattice is fixed, a tanh-type or elliptic ansatz with one pole per
period can return at most \(N_p\), respectively \(M_p\), normalized
solutions, independently of the expansion method.
This agrees with the critical comparisons of these methods
\citep{KudryashovLoguinova2009,Kudryashov2009Errors,KudryashovSoukharev2009}:
the ``new'' waves they returned for the KdV--Burgers and
Kuramoto--Sivashinsky equations coincide with the known ones
\citep{Kudryashov2009KdVB,KudryashovSoukharev2009}, and an ansatz whose
pole order is lower than \(p\) returns no wave at all
\citep{Kudryashov2010Kawahara}.
For the conformable derivative \(T_\alpha f(t)=t^{1-\alpha}f'(t)\),
\(t>0\), \(0<\alpha\le1\) \citep{KhalilEtAl2014}, the change
\(\tau=t^\alpha/\alpha\) gives \(T_\alpha=\dd/\dd\tau\), and similarly
for a spatial conformable derivative, so whenever the transformed
equation has the profile equation \eqref{eq:main} its waves fall under
the same counts; this does not apply to the Caputo or
Riemann--Liouville derivatives.

\section{Examples}\label{sec:examples}

The examples show the counts of Table~\ref{tab:named} at work and
connect them with the classical lists. Order three was begun in
Section~\ref{sec:ks-rational}.

\subsection{Order two: KdV--Burgers}\label{sec:kdvb}

The equation \(u_t+uu_x+\nu u_{xx}+\delta u_{xxx}=0\), \(\delta\ne0\),
has \(p=2\); the coefficient \(\nu\) may have either sign. It arises,
for instance, for waves in a liquid with gas bubbles
\citep{KudryashovSinelshchikovFourthOrder2010} and in a viscoelastic
tube \citep{KudryashovSinelshchikovChernyavsky2008}. Its exact front
solutions have been studied in several equivalent forms
\citep{Kudryashov1988,JeffreyXu1989,VliegHalford1991,Kudryashov2009Errors,Kudryashov2009KdVB}.
For \(A(\rho)=\rho+a\), polynomial division gives
\begin{equation}\label{eq:kdvb-matching}
 \remA S_A=-\frac{a(a-1)(a+1)}6,\qquad
 P_A=\rho^2+5a\rho+a^2-6a-1,\qquad
 v_A=12\bigl((a+1)Q-Q^2\bigr).
\end{equation}
Thus there are exactly \(N_2=3\) simple pairs, as
Theorem~\ref{thm:prime} states for \(2p-1=3\):
\begin{equation}\label{eq:kdvb-pairs}
 \begin{array}{c|c|c}
 a&P_A(\rho)&v_A\\ \hline
 -1&(\rho-2)(\rho-3)&-12Q^2\\
 0&\rho^2-1&12Q(1-Q)=3\operatorname{sech}^2(z/2)\\
 1&(\rho-1)(\rho+6)&12Q(2-Q).
 \end{array}
\end{equation}
The pulse \(a=0\) is the KdV soliton, the single pair of the purely
dispersive subfamily, \(F_2=1\). The two fronts belong to mixed
equations and are mirror images of each other: for a given equation
with \(\nu\ne0\) they are one and the same front, read in the two
directions (Proposition~\ref{prop:fixed}). In the original variables,
\eqref{eq:scaling} and \eqref{eq:coefficients} read
\[
 u(x,t)=w_\ast+\delta\kappa^2v_A\bigl(\kappa(x-ct)\bigr),\qquad
 \nu=\delta\kappa(5a),\qquad
 w_\ast-c=\delta\kappa^2(a^2-6a-1),
\]
so a nonzero \(\nu\) fixes the front rate to
\(\kappa=\nu/(5a\delta)\), \(a=\pm1\), in agreement with the classical
parameter restriction. The elliptic pairs at this order, both for the
KdV equation, are \eqref{eq:kdvb-elliptic}. Theorem~\ref{thm:complete}
applies to the KdV subfamily \(a_1=0\); for the mixed operators the
same conclusion cannot be expected (Remark~\ref{rem:fisher}).

\subsection{Order three: elliptic pairs, and the lattices on which they
merge}\label{sec:ks}

An elliptic solution must have zero residue at its pole. For
\(P=D^3+BD^2+\mu D+a_0\) the Laurent residue is
\(60\mu/19-15B^2/76\), so \(B^2=16\mu\); this is the curve
\(\sigma^2=16\) of \eqref{eq:ks-sigma}. With this restriction the
entire elliptic matching scheme is
\begin{equation}\label{eq:ks-elliptic-scheme}
 \begin{gathered}
 P=D^3+B D^2+\frac{B^2}{16}D+a_0,\qquad
 v=-60\wp'-15B\wp-\frac{B^3}{64}-a_0,\\
 B^4=3072g_2,\qquad
 a_0^2=\frac{13B^6}{2048}-2160g_3.
 \end{gathered}
\end{equation}
Read with the operator fixed, this is the classical solution of
\citet{Kudryashov1990}, see also \citet{Kudryashov1991}: \(w=v+a_0\)
solves the usual integrated equation
\(w'''+Bw''+\mu w'+w^2/2+K=0\) with \(K=-a_0^2/2\). For a given
equation with \(B^2=16\mu\), \(g_2\) is fixed by \(B\) and \(g_3\) is
free: a one-parameter family of elliptic waves. \citet{Hone2005} proves
that the restriction \(B^2=16\mu\) is necessary, and
\citet{Eremenko2006} completes the classification of all meromorphic KS
solutions. For the count we read \eqref{eq:ks-elliptic-scheme} the
other way: fix \(g_2,g_3\) and solve for \(B\) and \(a_0\).

On a generic lattice the equation for \(B\) has four roots and, for
each of them, the equation for \(a_0\) has two, the two backgrounds
\(\pm a_0\) of \eqref{eq:shift}: these are the \(M_3=8\) pairs. On
every lattice these two equations have eight solutions, counted with
multiplicity. This is the count of Theorem~\ref{thm:elliptic-count}:
the unknowns eliminated on the way to \eqref{eq:ks-elliptic-scheme} are
determined linearly, with nonzero rational pivots, so its last two
equations carry the full matching system, multiplicities included. On
special lattices pairs merge, and the multiplicity of a merged pair is
the number of pairs that have come together
(Section~\ref{sec:counting}). Pairs merge exactly where the Jacobian of
the two equations vanishes, and that Jacobian is a nonzero constant
times \(B^3a_0\). In terms of the invariant \(j(\Lambda)\) of
Section~\ref{sec:elliptic-counting}: if \(B=0\), then \(g_2=0\), hence
\(j=0\); if \(a_0=0\) on a lattice, elimination of \(B\) gives
\(675g_3^2=169g_2^3\), that is, \(j=-300\). These conditions are
sufficient as well. Thus:
\begin{equation}\label{eq:ks-collisions}
 \begin{array}{c|c|c}
 j(\Lambda)&\text{distinct pairs}&\text{total with multiplicity}\\ \hline
 j\notin\{0,-300\}&8&8\times1\\
 -300&6&2\times2+4\times1\\
 0&2&2\times4.
 \end{array}
\end{equation}
This is the exceptional set of Theorem~\ref{thm:elliptic-count} at
order three, and it shows the two ways in which pairs can meet. At
\(j=-300\) the four values of \(B\) remain distinct, and for two of
them the two backgrounds coincide at \(a_0=0\), in one pair of
multiplicity two; for instance, \(g_2=1/12\), \(g_3=13/1080\) gives the
double points \((B,a_0)=(\pm4,0)\). At \(j=0\) all four values of \(B\)
meet at \(B=0\), and the two roots of \(a_0^2=-2160g_3\) remain
distinct, since \(g_3\ne0\). Both pairs have the operator \(D^3+a_0\),
which is purely dissipative.

\subsection{Purely dispersive equations: Kawahara and seventh-order KdV}
\label{sec:kawahara}

\paragraph{Order four.}
For the Kawahara equation
\citep{Kawahara1972,DeminaKudryashovKawahara2010} the operator is
\begin{equation}\label{eq:kawahara-P}
 P=D^4+a_2D^2+a_0.
\end{equation}
The classification for this equation, using its first integral
\eqref{eq:energy}, is already established by
\citet{DeminaKudryashovKawahara2010}; see also \citet{YeEtAl2022}. By
Theorem~\ref{thm:symmetric}, the pairs with an even operator are those
with \(\iota A=A\), that is, \(A=\rho^3+\alpha_2\rho\). All of them
are pulses: the profile is even about its pole, \(v(Q)=v(1-Q)\). The
quotient and remainder are
\begin{equation}\label{eq:kawahara-exponential}
 \begin{split}
 P_A&=\rho^4+13\alpha_2\rho^2+\frac{31\alpha_2^2-70\alpha_2+7}{3},\\
 v_A&=280(1+\alpha_2)Q(1-Q)-1680Q^2(1-Q)^2,\\
 \remA S_A
   &=-\frac{\rho}{420}(\alpha_2+1)(31\alpha_2^2-31\alpha_2+10).
 \end{split}
\end{equation}
The three distinct roots, \(\alpha_2=-1\) and
\(\alpha_2=1/2\pm3i/(2\sqrt{31})\), give \(F_4=\binom31=3\) simple
pairs. The cubic already appears in \citet[Eq.~(47)]{Kudryashov1989},
with \(y=-1/\alpha_2\), where one real and two complex conjugate roots
are noted. The three values are the simply periodic parameter choices
of \citet[Section~3]{DeminaKudryashovKawahara2010}, whose parameter in
Eq.~(31) equals \(1/\alpha_2\). At \(\alpha_2=-1\),
\begin{equation}\label{eq:kawahara-soliton}
 v=-105\operatorname{sech}^4(z/2),\qquad
 P=(D^2-4)(D^2-9).
\end{equation}
This is the familiar solitary wave of the Kawahara equation; the other
two pairs are complex at unit real rate. Of all \(N_4=35\) pairs at
this order, \(21\) have real coefficients, and all \(21\) have rational
coefficients: five pulses and sixteen fronts, eleven wave shapes up to
reflection. They are the one-pole waves of the mixed fifth-order
equations, such as the Benney--Lin equation and the equation of
\citet{KudryashovSinelshchikov2011} for a viscoelastic tube, and the
repository dictionary lists them; solitary waves of this family are
found in \citet{KudryashovDemina2009}.

Substitution shows that the elliptic pairs of \eqref{eq:kawahara-P} are
given by
\begin{align}
 v&=-1680\wp^2-\frac{280a_2}{13}\wp
       +168g_2+\frac{31a_2^2}{507}-a_0,
       \label{eq:kawahara-elliptic}\\
 0&=31a_2^3-42588a_2g_2-4745520g_3,
       \label{eq:kawahara-cubic}\\
 a_0^2&=23184g_2^2+\frac{7980}{169}a_2^2g_2
                       +\frac{101520}{13}a_2g_3.
       \label{eq:kawahara-background}
\end{align}
These are the Weierstrass profiles and parameter relations obtained
from Laurent series by
\citet[Section~3]{DeminaKudryashovKawahara2010}; in their notation
\(a_2=-\beta\), \(a_0^2=C_0^2+12C_1\) and \(v=-6(w-w_\ast)\), where
their \(w_\ast\) is a constant equilibrium and \(a_0=C_0-6w_\ast\).
The first condition is cubic in \(a_2\), and the second is monic
quadratic in \(a_0\), so there are \(3\times2=6\) solutions, counted
with multiplicity, on every lattice: three values of \(a_2\), each with
its two backgrounds. This is the number \(2F_4\) of
Corollary~\ref{cor:elliptic-symmetric}. At \((g_2,g_3)=(1,0)\) the
cubic has three distinct roots and the right side of
\eqref{eq:kawahara-background} is nonzero at all of them, so the six
pairs are simple on a nonempty open set of lattices.
The unrestricted elliptic count at this order is \(M_4=30\); elliptic
formulas with nonzero odd derivative coefficients in the full
fifth-order evolution equation are given by
\citet{KudryashovSinelshchikov2012}.

\paragraph{Order six.}
The seventh-order KdV equation has the operator
\(P=D^6+a_4D^4+a_2D^2+a_0\); its meromorphic travelling waves are
studied by \citet{Dang2023Even}.
Here \(A=\rho^5+\alpha_2\rho^3+\alpha_4\rho\) and \(s_6=5544\). The
remainder of \(S_A\) has two coefficients, at \(\rho^3\) and at
\(\rho\), which are polynomials in \(\alpha_2,\alpha_4\) of weights
eight and ten. Eliminating \(\alpha_4\), which is determined by
\(\alpha_2\) at every solution, gives
\begin{equation}\label{eq:p6-eliminant}
 (\alpha_2+5)(\alpha_2+10)(49\alpha_2^2-5\alpha_2+10)\,
 \varphi_6(\alpha_2)=0,
\end{equation}
where \(\varphi_6\) is an irreducible sextic without real roots; the
ancillary script prints it and the two coefficients. The ten roots are
distinct, in agreement with \(F_6=\binom52=10\), and all ten pairs are
pulses. Two pairs are real, the solitary waves of
\citet{RyabovSinelshchikovKochanov2011}:
\begin{equation}\label{eq:p6-real}
 \begin{array}{c|c|c}
 A&P_A&v_A\\ \hline
 \rho(\rho^2-1)(\rho^2-4)&(\rho^2-9)(\rho^2-16)(\rho^2-25)
   &10395\operatorname{sech}^6(z/2)\\
 \rho(\rho^2-1)(\rho^2-9)&(\rho^2-4)(\rho^2-25)(\rho^2-71)
   &10395\operatorname{sech}^4(z/2)\bigl(1+\operatorname{sech}^2(z/2)\bigr).
 \end{array}
\end{equation}
An exact computation on a sample lattice gives \(2F_6=20\) for the
elliptic pairs of this family.

\subsection{Purely dissipative equations of order five}\label{sec:p5}

At \(p=5\) the purely dissipative operators are
\(P=D^5+a_3D^3+a_1D+a_0\), and \(A=\rho^4+\alpha_2\rho^2+\alpha_4\).
All \(126\) pairs at \(p=5\) are simple
(Appendix~\ref{app:verification}). By
Proposition~\ref{prop:general}(iii) there are therefore exactly
\(F_5=6\) distinct pairs here, and all six are fronts,
\(A(0)=\alpha_4\ne0\); for the four real ones this is also
Theorem~\ref{thm:symmetric}(iii). They are the three pairs in the
first line of
\begin{equation}\label{eq:p5-symmetric}
 \begin{gathered}
 A=(\rho^2-1)(\rho^2-4),\qquad
 A=(\rho^2-1)(\rho^2-9),\qquad
 A=(\rho^2-1)\Bigl(\rho^2+\frac15\Bigr),\\
 1105\alpha_2^3-2969\alpha_2^2+5570\alpha_2-4000=0,\qquad
 \alpha_4=\frac{5\alpha_2^3+21\alpha_2-10}{79\alpha_2+210},
 \end{gathered}
\end{equation}
and the three pairs given by the second line; the cubic has one real
root. The six pairs are the six solitary waves of
\citet{KudryashovDemina2007}, in whose notation \(w=1/(4a_3)\). The
four with real coefficients are, after the scales are restored, the
solutions~(1.4) of \citet{KudryashovMigita2007}; see also
\citet{RyabovSinelshchikovKochanov2011}. For instance,
\(A=(\rho^2-1)(\rho^2-4)\) gives
\(P_A=\rho^5-55\rho^3+1254\rho-2520\) and
\(v_A=5040\,Q^3(10-15Q+6Q^2)\), a front from \(0\) to
\(5040=-2P_A(0)\).

\section{Open questions and the boundary of the classification}\label{sec:open}

Two questions remain open in the enumeration, and one limitation of
the classification deserves a precise statement.

\paragraph{Rigidity.}
The first question is the continuous-convolution rigidity assertion
\[
 0\ne A\in\CC[\rho],\qquad
 A\mid\int_0^\rho A(t)A(\rho-t)\dd t
 \quad\Longrightarrow\quad A=c\rho^{\deg A}.
\]
Through the polynomial Laplace correspondence used in the proof of
Theorem~\ref{thm:elliptic-count}, it says that, after fixing the pole
at zero, any rational one-pole solution of \eqref{eq:main} vanishing at
infinity would have to be \(c_\ast z^{-p}\), with operator
\(P=D^p\).
Algebraic and \(\ell\)-adic arguments, together with exact
computations, prove it in the range
\(\deg A\le572\), corresponding to \(2\le p\le573\), and in infinite
arithmetic families of higher degrees.
Rigidity in arbitrary degree remains a conjecture.
By Proposition~\ref{prop:general}, it would establish both
enumeration formulas, as counts with multiplicity, at every order.

\paragraph{Simplicity.}
The second question is whether every point of \(\Xp\), the
rational-exponential matching system, is simple. The elliptic
simplicity assertion is generic in the lattice and is derived from
this one in Theorem~\ref{thm:elliptic-count}.
At a solution \(S_A=AC_A\), the tangent map of the system is
\begin{equation}\label{eq:tangent}
 J_A(h)=\remA\bigl(2(A\star h)-C_Ah\bigr),\qquad \deg h<d.
\end{equation}
The question is whether this map always has zero kernel.
The prime-index argument proves this after reduction modulo
\(\ell\), where the map becomes multiplication by a polynomial
coprime to \(\overline A\); a uniform replacement is missing at
composite indices.
One might expect simplicity to follow from the uniqueness of the
Laurent expansion in Theorem~\ref{thm:complete}. It does not: Laurent
uniqueness concerns a fixed operator, whereas \eqref{eq:tangent} allows
the operator to vary with the profile.

\begin{conjecture}\label{conj:count}
Continuous-convolution rigidity holds in every degree, and the
matching Jacobian \eqref{eq:tangent} is invertible at every point
of \(\Xp\). Consequently there are exactly \(\binom{2p-1}{p-1}\)
rational-exponential pairs, all distinct, for every \(p\ge2\).
\end{conjecture}

The conjecture would also imply \(M_p\) simple elliptic pairs
on a generic fixed lattice at every order, \(F_p\) distinct pairs
in the two pure subfamilies, and that the rational-exponential waves of
every purely dissipative equation are fronts
(Theorem~\ref{thm:symmetric}).
Simplicity on every lattice is false already at order two;
\eqref{eq:ks-collisions} describes the full exceptional set
at order three.

\paragraph{Mixed equations with \(p\) even.}
The hypothesis of Theorem~\ref{thm:complete} is a separate issue. For
even \(p\) it requires a purely dispersive equation, and it cannot
simply be dropped. The known example comes from the Fisher equation.
Its travelling waves of nonzero speed satisfy \eqref{eq:main} with
\(p=2\) and a mixed operator, which is also the profile equation of
KdV--Burgers. At the
special speed of \citet{AblowitzZeppetella1979} it has a one-parameter
family of solutions, given in \citet{Kudryashov2010},
\citet[Lemma~4.4]{ConteNgWu2015} and
\citet{McCueElHachemSimpson2021}; see also
\citet{KudryashovFisher2005}. In our normalization it is
\[
 v(z)=-12\lambda^2e^{2\lambda z}
       \wp(e^{\lambda z}-1;0,g_3),\qquad
 P(D)=D^2-5\lambda D+6\lambda^2,\qquad\lambda\ne0.
\]
Substitution follows from
\((D^2-5\lambda D+6\lambda^2)[t^2U(t)]
=\lambda^2t^4U''(t)\), \(t=e^{\lambda z}\).
For \(g_3=0\) and \(\lambda=1\) this is the front in the first row of
\eqref{eq:kdvb-pairs}. For generic \(g_3\ne0\) it is globally
meromorphic and lies outside \(\Wclass\).
Further examples of meromorphic solutions outside \(\Wclass\)
are discussed by \citet[Remark~1.11]{NgWu2019}.
Thus a classification for mixed operators with \(p\) even must allow a
larger class of functions than \(\Wclass\).

For mixed operators of even degree \(p\ge4\) we do not know whether
meromorphic solutions outside \(\Wclass\) occur.
Methods for constructing elliptic and rational-exponential
solutions without the finiteness property are developed by
\citet{Demina2022}.

\appendix
\section{Exact verification and ancillary material}\label{app:verification}

The five kinds of computer checks below use exact rational or
finite-field arithmetic; instructions, dependency versions and checksums
are included with the source package. The prime-index theorem does not
depend on the computer calculations.
\begin{enumerate}\setlength{\itemsep}{0pt}\setlength{\parskip}{0pt}
\item \texttt{verify\_manuscript.py} reconstructs the convolutions, the
first Laurent coefficients and the Weierstrass relations, and checks the
leading-term comparison in Proposition~\ref{prop:injective} and, by
exact substitution, the displayed formulas of orders two to four,
including the agreement of Table~\ref{tab:ks} with the tables of
\citet{KudryashovZargaryan1996,ConteMusette2009}, the multiplicities
\eqref{eq:ks-collisions}, the Kawahara coefficient restriction and its
literature normalization, the nodal correspondence and the Fisher
identity. The uniform proofs are in the text.
\item \texttt{p5\_count.sing} (\textsc{Singular}, Gr\"obner basis over
\(\QQ\)): at \(p=5\) the count with multiplicity of the full matching
equations, their quotient dimension, is \(126\). Adjoining the Jacobian
determinant gives the unit ideal: the determinant vanishes at no
solution, so all are simple.
\item At \(p=8\) the same two outputs, \(6435\) and the unit ideal, are
obtained over the finite field \(\FF_{32003}\). So every solution modulo
\(32003\) is simple and lifts uniquely by Hensel's lemma, as in Step~4
of the proof of Theorem~\ref{thm:prime}, and the lifts are all the
solutions because the count with multiplicity is \(6435\)
(Proposition~\ref{prop:general}(ii)). At \(p=11\) each of the \(352716\)
solutions is certified over the \(19\)-adic numbers, by Hensel's lemma
at the simple points modulo \(19\) and by exact Newton-polygon and
Newton--Kantorovich certificates at the singular ones. The archived
certificate and the record of its validation are in the repository
(\texttt{certificates/}).
\item \texttt{rigidity/} has the proof and exact certificate for
rigidity \((\mathrm R_d)\), \(d=p-1\le72\). For \(2d+1=\ell^e+g\),
\(g\in\{0,2,4,6\}\), the proof reduces the normalized continuous
convolution to a residual system of degree at most \(g\); the
prime-power case \(g=0\) is also in the companion note. The certificate
retains every residual degree, with integer ideal-membership identities
and certified prime factors; the replay multiplies out every identity
and checks all 72 arithmetic choices. It does not infer
characteristic-zero nonexistence from lists of residue-field points.
\item \texttt{symmetry\_checks/} checks Proposition~\ref{prop:reflection}
with symbolic coefficients for \(d\le5\); the labels of
Table~\ref{tab:ks} and of the \(21\) real pairs at \(p=4\), with
Proposition~\ref{cor:endpoints} and Theorem~\ref{thm:symmetric}(ii) for
them; the counts with multiplicity \(F_p\) for \(p\le8\), and
finite-field certificates, as at \(p=8\), that the \(F_p\) pairs are
distinct for every \(p\le14\); on a sample lattice, the counts with
multiplicity \(2F_p\) of Corollary~\ref{cor:elliptic-symmetric} for
\(p\le6\), and no elliptic pair of a purely dissipative equation for
\(p=3,5,7\) (\texttt{elliptic\_dissipative.py}); the six pairs
\eqref{eq:p5-symmetric} and the ten pairs \eqref{eq:p6-eliminant}, with
the sextic factor in full; and the scaling \eqref{eq:fixed-wave} by
substitution in the evolution equation.
\end{enumerate}

The Lean~4 formalization of Remark~\ref{rem:lean} is at
\url{https://github.com/migita/binomial-count-meromorphic-waves},
directory \texttt{lean/}, with Lean and Mathlib pinned to
version~4.33.1. It proves rigidity \((\mathrm R_d)\) for \(2d+1\) prime
(\texttt{rigidity\_prime\_index}) and the count and simplicity
statements of Theorem~\ref{thm:prime}, both for the matching equations
\eqref{eq:remainder} (\texttt{primeIndexCountAndSimplicity}) and for
the normalized pairs \((P,v)\) as solutions of \eqref{eq:main} in the
field of rational functions of \(e^{z}\)
(\texttt{primeIndexRationalExponentialCountAndSimplicity}), with the
correspondence between the two formulations for every \(p\ge2\). The
discrete convolution is defined there by the Bernoulli-polynomial power
sums; the remainder equations and their Jacobian are those of
Section~\ref{sec:counting}. \texttt{EnumerationCheck.lean} and
\texttt{Check.lean} verify that the proofs have exactly the stated
propositions as their types and depend only on the axioms
\texttt{propext}, \texttt{Classical.choice} and \texttt{Quot.sound}; the
project contains no \texttt{sorry}, added axiom or
\texttt{native\_decide}.


\section*{Tool and computational resource disclosure}

This paper was written with substantial use of generative AI, which is
disclosed here in the form recommended by the Leiden Declaration on
Artificial Intelligence and
Mathematics\footnote{\url{https://leidendeclaration.ai/}, 2 June 2026.}
and required by arXiv.
No automated system is an author; credit and responsibility for the
results, and for the attribution of prior work, rest with the author.

\begin{itemize}
\item \texttt{gpt-6-astra} (OpenAI), August--September 2026. The main
mathematical ideas of the paper, the proofs of
Theorems~\ref{thm:complete}, \ref{thm:prime}
and~\ref{thm:elliptic-count} and of the weighted count in the companion
note, the rigidity condition
\((\mathrm R_d)\) and its certificate, the exact verification scripts,
the \textsc{Singular} input, the first complete draft of the text,
and the Lean~4 formalization of Remark~\ref{rem:lean} were produced by
this model in dialogue with the author.
\item OpenAI Codex v0.153.4, September 2026. Independent mathematical
audit of the final revision, integration of the edits, and rebuilding
and validation of the source package.
\item Claude Code (Anthropic, Claude Fable~5.1), September 2026. The
literature search behind the attributions in Section~\ref{sec:intro},
the readability revision of all sections, and the leading-coefficient
argument of Proposition~\ref{prop:injective}, which Codex then checked
independently. In the final revision it proposed and proved the
results on the two pure subfamilies and on single equations
(Corollary~\ref{cor:symmetric-waves}, Section~\ref{sec:reflection},
Corollary~\ref{cor:elliptic-symmetric} and Section~\ref{sec:fixed}),
with their exact checks, and reorganized the manuscript around them.
In the last revision it made the prime-index theorem the spine of the
paper: it wrote the proof of Theorem~\ref{thm:prime} in the form used
by the Lean formalization, derived
Proposition~\ref{cor:endpoints} and Theorem~\ref{thm:symmetric} from
the labels, replaced the B\'ezout count in the proof of
Theorem~\ref{thm:elliptic-count} by a flat-family argument, and moved
the general-order material to the companion note; Codex reviewed these
arguments. It also audited the Lean formalization independently:
the formal statements against the manuscript, a rebuild from source,
and the axiom audit.
\end{itemize}

{\small\raggedright
\setlength{\bibsep}{3pt}
\interlinepenalty=10000
\bibliographystyle{plainnat}
\bibliography{references}
}
\end{document}
